\paragraph{属 \(49\) 的 \texorpdfstring{$S_4$}{S4}--Hurwitz 商塔上的商曲线亏格、总分歧预算、Prym Weil 型、局部 \texorpdfstring{$\zeta$}{zeta} 分解与 Hodge 类传递}
沿用定理 \ref{thm:killo-s4-even-subgroup-free-action-unramified-a4-lift}、\ref{thm:killo-s4-burnside-kani-rosen-prym-square}、\ref{thm:killo-s4-genus49-jacobian-computable-isogeny} 与命题 \ref{prop:killo-two-prym-polarization-types} 的记号，记
\[
X:=\mathcal X,\qquad
H:=\mathcal H=X/A_4,\qquad
C_6:=\mathcal C_6=X/V_4,
\]
\[
Z:=X/D_8,\qquad
C_F:=\mathcal C_F=X/S_3,\qquad
Y:=X/C_4,
\]
并置
\[
A:=\operatorname{Prym}(C_6/H),\qquad
P:=\operatorname{Prym}(Y/Z).
\]

在这一记号下，商塔各层亏格、总分歧预算、分支纤维的固定点几何、Chevalley--Weil 重数的可逆恢复与两类 Prym 极化缺陷的素数分裂都可进一步完全显式化。

\begin{conclusion}[属 \(49\) 的 \texorpdfstring{$S_4$}{S4}--Hurwitz 曲线上四类非平凡循环商的亏格光谱]\label{con:xi-time-part9n-s4-cyclic-quotient-genus-spectrum}
设
\[
V:=H^0(X,\Omega_X^1)
\cong
5\cdot \mathrm{sgn}\ \oplus\ 4\cdot V_2\ \oplus\ 3\cdot V_3\ \oplus\ 9\cdot V_3',
\]
其中平凡表示不出现。若 \(\tau,\delta,\sigma_3,\sigma_4\in S_4\) 分别表示换位、双换位、\(3\)-循环与 \(4\)-循环，则
\[
g(X/\langle \tau\rangle)=19,\qquad
g(X/\langle \delta\rangle)=25,\qquad
g(X/\langle \sigma_3\rangle)=17,\qquad
g(X/\langle \sigma_4\rangle)=13.
\]
\end{conclusion}
\begin{proof}
对任意子群 \(U\le S_4\)，恒有
\[
g(X/U)=\dim V^U.
\]
因此只需计算 \(V\) 在四类非平凡循环子群下的不变维数。取 \(S_4\) 的角色值
\[
\begin{array}{c|cccc}
 & \tau & \delta & \sigma_3 & \sigma_4\\
\hline
\mathrm{sgn} & -1 & 1 & 1 & -1\\
V_2 & 0 & 2 & -1 & 0\\
V_3 & 1 & -1 & 0 & -1\\
V_3' & -1 & -1 & 0 & 1
\end{array}
\]
并用角色平均公式
\[
\dim W^{\langle \tau\rangle}=\frac{\chi_W(1)+\chi_W(\tau)}{2},\qquad
\dim W^{\langle \delta\rangle}=\frac{\chi_W(1)+\chi_W(\delta)}{2},
\]
\[
\dim W^{\langle \sigma_3\rangle}=\frac{\chi_W(1)+2\chi_W(\sigma_3)}{3},\qquad
\dim W^{\langle \sigma_4\rangle}=\frac{\chi_W(1)+\chi_W(\sigma_4)+\chi_W(\delta)+\chi_W(\sigma_4^{-1})}{4},
\]
可得
\[
\dim(\mathrm{sgn}^{\langle\tau\rangle},V_2^{\langle\tau\rangle},V_3^{\langle\tau\rangle},(V_3')^{\langle\tau\rangle})
=(0,1,2,1),
\]
\[
\dim(\mathrm{sgn}^{\langle\delta\rangle},V_2^{\langle\delta\rangle},V_3^{\langle\delta\rangle},(V_3')^{\langle\delta\rangle})
=(1,2,1,1),
\]
\[
\dim(\mathrm{sgn}^{\langle\sigma_3\rangle},V_2^{\langle\sigma_3\rangle},V_3^{\langle\sigma_3\rangle},(V_3')^{\langle\sigma_3\rangle})
=(1,0,1,1),
\]
\[
\dim(\mathrm{sgn}^{\langle\sigma_4\rangle},V_2^{\langle\sigma_4\rangle},V_3^{\langle\sigma_4\rangle},(V_3')^{\langle\sigma_4\rangle})
=(0,1,0,1).
\]
于是分别得到
\[
g(X/\langle \tau\rangle)=5\cdot0+4\cdot1+3\cdot2+9\cdot1=19,
\]
\[
g(X/\langle \delta\rangle)=5\cdot1+4\cdot2+3\cdot1+9\cdot1=25,
\]
\[
g(X/\langle \sigma_3\rangle)=5\cdot1+4\cdot0+3\cdot1+9\cdot1=17,
\]
\[
g(X/\langle \sigma_4\rangle)=5\cdot0+4\cdot1+3\cdot0+9\cdot1=13.
\]
证毕。
\end{proof}

\begin{conclusion}[分支纤维的六对换位固定点分解与其余三类循环元的自由作用]\label{con:xi-time-part9n-s4-branch-fiber-transposition-fixedpairs}
设
\[
f:X\longrightarrow \PP^1_u
\]
为属 \(49\) 的 \(S_4\)-正则闭包覆盖。由结论 \ref{con:fold-gauge-anomaly-s4-hurwitz-tower-genus}，其分支集合为
\[
\{u=0,1\}\cup\{P_{10}(u)=0\},
\]
共 \(12\) 个分支值，且每个分支点的惯性群都由换位生成。于是：
\begin{enumerate}
  \item 对任意分支值 \(u_0\)，纤维 \(f^{-1}(u_0)\) 恰有 \(12\) 个点；
  \item 固定任一换位 \(\tau\in S_4\)。则在每个分支纤维上，\(\tau\) 恰固定 \(2\) 个点；因此 \(\tau\) 在整条曲线上恰固定 \(24\) 个点；
  \item 每个分支纤维因而可分裂为六对两两不交的换位固定点，每一对对应于 \(S_4\) 的一条换位；
  \item 任一双换位、\(3\)-循环或 \(4\)-循环在 \(X\) 上都无固定点。
\end{enumerate}
\end{conclusion}
\begin{proof}
取任一分支值 \(u_0\)，其惯性群记为
\[
I_{u_0}=\langle \tau_0\rangle\simeq C_2.
\]
由于 \(f\) 为正则 \(S_4\)-覆盖，有标准陪集识别
\[
f^{-1}(u_0)\cong S_4/I_{u_0},
\]
故
\[
\#f^{-1}(u_0)=|S_4:I_{u_0}|=24/2=12.
\]

现固定一条换位 \(\tau\)。对任意分支纤维，由正则覆盖的陪集公式，
\[
\#\bigl(\operatorname{Fix}_X(\tau)\cap f^{-1}(u_0)\bigr)
=
\frac{\#\{x\in S_4:\ x^{-1}\tau x\in I_{u_0}\}}{|I_{u_0}|}.
\]
而 \(I_{u_0}\) 的唯一非平凡元即 \(\tau_0\)，故上式条件等价于
\[
x^{-1}\tau x=\tau_0.
\]
满足此式的 \(x\) 构成一个 \(C_{S_4}(\tau)\)-左陪集；又 \(|C_{S_4}(\tau)|=4\)，从而
\[
\#\bigl(\operatorname{Fix}_X(\tau)\cap f^{-1}(u_0)\bigr)=4/2=2.
\]
对全部 \(12\) 个分支值求和即得
\[
\#\operatorname{Fix}_X(\tau)=12\cdot 2=24.
\]
另一方面，每个分支纤维中任一点的稳定子都等于某个换位子群 \(xI_{u_0}x^{-1}\)，故该点只可能被唯一的非平凡换位固定。由 \(6\) 条换位各贡献 \(2\) 个固定点，整个分支纤维遂分解成六对彼此不交的换位固定点。

最后，对双换位 \(\delta\)、\(3\)-循环 \(\sigma_3\) 与 \(4\)-循环 \(\sigma_4\)，由结论 \ref{con:xi-time-part9n-s4-cyclic-quotient-genus-spectrum} 及 Riemann--Hurwitz 公式分别有
\[
96=2\bigl(2g(X/\langle\delta\rangle)-2\bigr)+R_\delta
=2(2\cdot 25-2)+R_\delta,
\]
\[
96=3\bigl(2g(X/\langle\sigma_3\rangle)-2\bigr)+R_{\sigma_3}
=3(2\cdot 17-2)+R_{\sigma_3},
\]
\[
96=4\bigl(2g(X/\langle\sigma_4\rangle)-2\bigr)+R_{\sigma_4}
=4(2\cdot 13-2)+R_{\sigma_4},
\]
故 \(R_\delta=R_{\sigma_3}=R_{\sigma_4}=0\)。因此这三类循环元都自由作用于 \(X\)。证毕。
\end{proof}

\begin{conclusion}[五个循环商亏格对 Chevalley--Weil 重数向量的线性刚性]\label{con:xi-time-part9n-s4-quotient-genera-recover-multiplicities}
设
\[
H^0(X,\Omega_X^1)\cong
m_{\mathbf 1}\cdot \mathbf 1\oplus
m_{\mathrm{sgn}}\cdot \mathrm{sgn}\oplus
m_2\cdot V_2\oplus
m_3\cdot V_3\oplus
m_3'\cdot V_3'.
\]
则仅由
\[
g(X/S_4)=0,\qquad
g(X/\langle \tau\rangle)=19,\qquad
g(X/\langle \delta\rangle)=25,
\]
\[
g(X/\langle \sigma_3\rangle)=17,\qquad
g(X/\langle \sigma_4\rangle)=13
\]
这五个 genus 数据，便可唯一恢复
\[
(m_{\mathbf 1},m_{\mathrm{sgn}},m_2,m_3,m_3')=(0,5,4,3,9).
\]
\end{conclusion}
\begin{proof}
由 \(g(X/S_4)=\dim H^0(X,\Omega_X^1)^{S_4}\) 立即得到
\[
m_{\mathbf 1}=0.
\]
再由结论 \ref{con:xi-time-part9n-s4-cyclic-quotient-genus-spectrum} 的四类循环商 genus 与相应不变维数公式，得到线性系统
\[
m_2+2m_3+m_3'=19,
\]
\[
m_{\mathrm{sgn}}+2m_2+m_3+m_3'=25,
\]
\[
m_{\mathrm{sgn}}+m_3+m_3'=17,
\]
\[
m_2+m_3'=13.
\]
由第一式与第四式相减，得
\[
2m_3=6,\qquad m_3=3.
\]
代回第三式得
\[
m_{\mathrm{sgn}}+m_3'=14.
\]
再将第四式代入第二式，可得
\[
m_{\mathrm{sgn}}+m_2=12.
\]
与
\[
m_{\mathrm{sgn}}+m_3'=14,\qquad m_2+m_3'=13
\]
联立，推出
\[
m_2=4,\qquad m_3'=9,\qquad m_{\mathrm{sgn}}=5.
\]
于是得到唯一解
\[
(m_{\mathbf 1},m_{\mathrm{sgn}},m_2,m_3,m_3')=(0,5,4,3,9).
\]
证毕。
\end{proof}

\begin{conclusion}[八个商曲线的属谱完全表]\label{con:xi-time-part9n-s4-all-subgroup-genera}
设
\[
V:=H^0(X,\Omega_X^1)
\cong 5\cdot \mathrm{sgn}\oplus 4\cdot V_2\oplus 3\cdot V_3\oplus 9\cdot V_3'
\]
为属 \(49\) 的 \texorpdfstring{$S_4$}{S4}--Hurwitz 曲线 \(X\) 的微分表示。记
\[
H:=X/A_4,\qquad
C_6:=X/V_4,\qquad
Z:=X/D_8,\qquad
C_F:=X/S_3,\qquad
Y:=X/C_4,
\]
并记
\[
W_3:=X/C_3,\qquad W_2:=X/C_2.
\]
其中 \(C_2=\langle\tau\rangle\)、\(C_3=\langle\sigma_3\rangle\)、\(C_4=\langle\sigma_4\rangle\) 分别表示由换位、\(3\)-循环与 \(4\)-循环生成的循环子群。则所有这些商曲线的亏格都由 \(S_4\)-微分表示唯一决定，并满足
\[
g(X/S_4)=0,\qquad
g(H)=g(X/A_4)=5,\qquad
g(C_6)=g(X/V_4)=13,\qquad
g(Z)=g(X/D_8)=4,
\]
\[
g(C_F)=g(X/S_3)=3,\qquad
g(Y)=g(X/C_4)=13,\qquad
g(W_3)=g(X/C_3)=17,\qquad
g(W_2)=g(X/C_2)=19.
\]
特别地，\(W_3\) 与 \(W_2\) 的亏格无需额外 Hurwitz 计算，已由同一组 Chevalley--Weil 重数完全锁定。
\end{conclusion}
\begin{proof}
对任意子群 \(K\le S_4\)，商曲线亏格满足
\[
g(X/K)=\dim V^K.
\]
记
\[
\rho_K:=\operatorname{Ind}_{K}^{S_4}\mathbf 1.
\]
由 Frobenius 互反，
\[
\dim(\pi^K)=\langle \pi,\rho_K\rangle_{S_4}
\]
对每个不可约表示 \(\pi\) 都成立。另一方面，定理 \ref{thm:killo-s4-burnside-kani-rosen-prym-square} 给出
\[
\rho_{A_4}=\mathbf 1\oplus \mathrm{sgn},\qquad
\rho_{V_4}=\mathbf 1\oplus \mathrm{sgn}\oplus 2V_2,\qquad
\rho_{D_8}=\mathbf 1\oplus V_2,
\]
\[
\rho_{S_3}=\mathbf 1\oplus V_3,\qquad
\rho_{C_4}=\mathbf 1\oplus V_2\oplus V_3',\qquad
\rho_{C_2}=\mathbf 1\oplus V_2\oplus 2V_3\oplus V_3',
\]
\[
\rho_{C_3}=\mathbf 1\oplus \mathrm{sgn}\oplus V_3\oplus V_3',\qquad
\rho_{S_4}=\mathbf 1.
\]
将这些分解与 \(V\) 的 Chevalley--Weil 分解逐项配对，便得到
\[
g(X/A_4)=5,\quad
g(X/V_4)=5+2\cdot4=13,\quad
g(X/D_8)=4,\quad
g(X/S_3)=3,
\]
\[
g(X/C_4)=4+9=13,\quad
g(X/C_3)=5+3+9=17,\quad
g(X/C_2)=4+2\cdot3+9=19.
\]
又因 \(V\) 中不含平凡表示，故
\[
g(X/S_4)=\dim V^{S_4}=0.
\]
证毕。
\end{proof}

\begin{conclusion}[属 \texorpdfstring{$49$}{49} 的 \texorpdfstring{$S_4$}{S4}--Hurwitz 链中间商映射总分歧度的 \texorpdfstring{$24$}{24} 倍数谱]\label{con:xi-time-part9n-s4-total-ramification-spectrum}
对上述各子群 \(K\le S_4\)，记
\[
\pi_K:X\to X/K
\]
为自然商映射，并令
\[
R_K:=\sum_{P\in X}\bigl(e_P(\pi_K)-1\bigr)
\]
为其总分歧度。则有显式公式
\[
R_{S_4}=144,\qquad
R_{A_4}=0,\qquad
R_{V_4}=0,\qquad
R_{D_8}=48,
\]
\[
R_{S_3}=72,\qquad
R_{C_4}=0,\qquad
R_{C_3}=0,\qquad
R_{C_2}=24.
\]
因此全部非零总分歧度恰落在
\[
24\cdot\{1,2,3,6\},
\]
并形成严格链
\[
24<48<72<144.
\]
\end{conclusion}
\begin{proof}
由上一定理与 Chevalley--Weil 分解可知 \(g(X)=49\)，故
\[
2g(X)-2=96.
\]
对任意 \(K\le S_4\)，Riemann--Hurwitz 公式给出
\[
2g(X)-2=|K|\bigl(2g(X/K)-2\bigr)+R_K,
\]
即
\[
R_K=96-|K|\bigl(2g(X/K)-2\bigr).
\]
将结论 \ref{con:xi-time-part9n-s4-all-subgroup-genera} 的亏格值与
\[
|S_4|=24,\quad |A_4|=12,\quad |V_4|=4,\quad |D_8|=8,\quad
|S_3|=6,\quad |C_4|=4,\quad |C_3|=3,\quad |C_2|=2
\]
逐一代入，立即得到
\[
R_{S_4}=144,\quad R_{A_4}=0,\quad R_{V_4}=0,\quad R_{D_8}=48,
\]
\[
R_{S_3}=72,\quad R_{C_4}=0,\quad R_{C_3}=0,\quad R_{C_2}=24.
\]
其中 \(R_{A_4}=R_{V_4}=R_{C_3}=0\) 与换位惯性对子群 \(A_4,V_4,C_3\) 的平凡交相一致，而 \(R_{C_4}=0\) 则对应于 \(C_4\) 中不含任何换位惯性元。证毕。
\end{proof}

\begin{conclusion}[属 \texorpdfstring{$49$}{49} 的 \texorpdfstring{$S_4$}{S4}--Hurwitz Jacobian 的 \texorpdfstring{$13+9+9+9+9$}{13+9+9+9+9} 型块压缩]\label{con:xi-time-part9n-jacobian-c6-compression}
设
\[
J(X)\sim_{\QQ}J(H)\times J(Z)^2\times J(C_F)^3\times P^3,
\qquad
P=\operatorname{Prym}(Y/Z),
\]
并沿上文记号令 \(C_6=X/V_4\)。则有更压缩的同源分解
\[
\boxed{\;
J(X)\sim_{\QQ} J(C_6)\times J(C_F)^3\times P^3.
\;}
\]
等价地，原来的
\[
5+2\cdot 4+3\cdot 3+3\cdot 9=49
\]
维分解可压缩为
\[
13+9+9+9+9=49
\]
的块结构，其中 \(13\) 维块 \(J(C_6)\) 已吸收 \(J(H)\) 与 \(J(Z)^2\) 的全部几何信息。
\end{conclusion}
\begin{proof}
定理 \ref{thm:killo-s4-burnside-kani-rosen-prym-square} 给出
\[
J(C_6)\sim_{\QQ}J(H)\times J(Z)^2,
\]
其几何内容正是无分歧循环三重覆盖 \(C_6\to H\) 的 Prym 分裂
\[
\operatorname{Prym}(C_6/H)\sim_{\QQ}J(Z)^2.
\]
另一方面，定理 \ref{thm:killo-s4-genus49-jacobian-computable-isogeny} 给出
\[
J(X)\sim_{\QQ}J(H)\times J(Z)^2\times J(C_F)^3\times P^3.
\]
将前一式代入后一式，立即得到
\[
J(X)\sim_{\QQ}J(C_6)\times J(C_F)^3\times P^3.
\]
又由
\[
g(C_6)=13,\qquad 3g(C_F)=9,\qquad \dim P^3=27=9+9+9
\]
可知该分解正对应于 \(13+9+9+9+9\) 的块压缩。证毕。
\end{proof}

\begin{theorem}[属 \texorpdfstring{$49$}{49} 的 \texorpdfstring{$S_4$}{S4}--Hurwitz Jacobian 具有显式半单非交换自同态子代数]\label{thm:xi-time-part9n-jacobian-endomorphism-lower-bound}
\leanverified{paper\_xi\_time\_part9n\_jacobian\_endomorphism\_lower\_bound}
沿用结论 \ref{con:xi-time-part9n-jacobian-c6-compression} 的记号，并令
\[
A:=J(H),\qquad
B:=J(Z),\qquad
C:=J(C_F),\qquad
D:=P.
\]
则有理自同态代数满足显式包含
\[
\boxed{\;
\operatorname{End}^0_{\QQ}\!\bigl(J(X)\bigr)
\supseteq
\QQ\times M_2(\QQ)\times M_3(\QQ)\times M_3(\QQ).\;}
\]
特别地，
\[
\boxed{\;
\dim_{\QQ}\operatorname{End}^0_{\QQ}\!\bigl(J(X)\bigr)\ge 1+2^2+3^2+3^2=23.\;}
\]
并且在 \(\QQ\)-同源意义下，\(J(X)\) 既不简单，也不等型。
\end{theorem}
\begin{proof}
由定理 \ref{thm:killo-s4-genus49-jacobian-computable-isogeny}，
\[
J(X)\sim_{\QQ}A\times B^2\times C^3\times D^3.
\]
同源不改变有理自同态代数，故
\[
\operatorname{End}^0_{\QQ}\!\bigl(J(X)\bigr)
\cong
\operatorname{End}^0_{\QQ}\!\bigl(A\times B^2\times C^3\times D^3\bigr).
\]
对任意阿贝尔簇 \(E\) 与整数 \(n\ge 1\)，标量矩阵给出自然嵌入
\[
M_n(\QQ)\hookrightarrow M_n\!\bigl(\operatorname{End}^0_{\QQ}(E)\bigr)
\cong
\operatorname{End}^0_{\QQ}(E^n).
\]
分别取 \(n=1,2,3,3\) 并作用于四个直积块，即得块对角嵌入
\[
\QQ\times M_2(\QQ)\times M_3(\QQ)\times M_3(\QQ)
\hookrightarrow
\operatorname{End}^0_{\QQ}\!\bigl(J(X)\bigr).
\]
于是
\[
\dim_{\QQ}\operatorname{End}^0_{\QQ}\!\bigl(J(X)\bigr)\ge 1+4+9+9=23.
\]

另一方面，四个可见因子的维数分别为
\[
\dim A=5,\qquad \dim B=4,\qquad \dim C=3,\qquad \dim D=9.
\]
不同维数的阿贝尔簇不可能互相同源，故 \(A,B,C,D\) 代表四个互不同源的等型块。因而 \(J(X)\) 在 \(\QQ\)-同源意义下至少分裂为四个互异块，既不简单，也不等型。证毕。
\end{proof}

\begin{theorem}[\texorpdfstring{$S_4$}{S4}--Hurwitz 链的圆维预算完全闭合]\label{thm:xi-time-part9n-s4-hurwitz-circle-dimension-budget}
\leanverified{paper\_xi\_time\_part9n\_s4\_hurwitz\_circle\_dimension\_budget}
沿用结论 \ref{con:xi-time-part9n-s4-all-subgroup-genera} 与
\ref{con:xi-time-part9n-jacobian-c6-compression} 的记号。对任意复光滑射影曲线 \(C\) 与复阿贝尔簇 \(A\)，记
\[
\Lambda(C):=H_1\!\bigl(C(\CC),\ZZ\bigr),
\qquad
\Lambda(A):=H_1\!\bigl(A(\CC),\ZZ\bigr).
\]
则
\[
\cdim\!\bigl(\Lambda(H)\bigr)=10,\qquad
\cdim\!\bigl(\Lambda(Z)\bigr)=8,\qquad
\cdim\!\bigl(\Lambda(C_F)\bigr)=6,\qquad
\cdim\!\bigl(\Lambda(P)\bigr)=18,
\]
\[
\cdim\!\bigl(\Lambda(C_6)\bigr)=26,
\qquad
\cdim\!\bigl(\Lambda(Y)\bigr)=26.
\]
并且整同调圆维预算满足严格闭合公式
\[
\boxed{\;
\cdim\!\bigl(\Lambda(X)\bigr)
=
\cdim\!\bigl(\Lambda(H)\bigr)
+2\,\cdim\!\bigl(\Lambda(Z)\bigr)
+3\,\cdim\!\bigl(\Lambda(C_F)\bigr)
+3\,\cdim\!\bigl(\Lambda(P)\bigr)
=98.\;}
\]
等价地，
\[
\boxed{\;
\cdim\!\bigl(\Lambda(X)\bigr)=98=10+2\cdot 8+3\cdot 6+3\cdot 18.\;}
\]
因此，属 \(49\) 的 Jacobian 在圆维层面不存在任何隐藏自由秩；文中给出的四块分解
\[
J(H),\qquad J(Z)^2,\qquad J(C_F)^3,\qquad P^3
\]
恰好耗尽全部 \(98\) 维整同调预算。
\end{theorem}
\begin{proof}
由定理 \ref{thm:killo-cdim-short-exact-additivity}，对任意有限生成离散交换群 \(G\) 有
\[
\cdim(G)=\operatorname{rk}_{\ZZ}(G)=\dim_{\QQ}(G\otimes_{\ZZ}\QQ),
\]
并且短正合列下圆维可加。

对任意亏格为 \(g(C)\) 的复光滑射影曲线 \(C\)，其一维整同调群是秩 \(2g(C)\) 的自由阿贝尔群，故
\[
\cdim\!\bigl(\Lambda(C)\bigr)=2g(C).
\]
对任意维数为 \(d\) 的复阿贝尔簇 \(A\)，同理有
\[
\cdim\!\bigl(\Lambda(A)\bigr)=2d.
\]
由结论 \ref{con:xi-time-part9n-s4-all-subgroup-genera}，
\[
g(H)=5,\qquad g(Z)=4,\qquad g(C_F)=3,\qquad g(C_6)=13,\qquad g(Y)=13.
\]
再由结论 \ref{con:xi-time-part9n-jacobian-c6-compression} 的记号，\(\dim P=9\)。于是
\[
\cdim\!\bigl(\Lambda(H)\bigr)=10,\qquad
\cdim\!\bigl(\Lambda(Z)\bigr)=8,\qquad
\cdim\!\bigl(\Lambda(C_F)\bigr)=6,
\]
\[
\cdim\!\bigl(\Lambda(P)\bigr)=18,\qquad
\cdim\!\bigl(\Lambda(C_6)\bigr)=26,\qquad
\cdim\!\bigl(\Lambda(Y)\bigr)=26.
\]

另一方面，由定理 \ref{thm:killo-s4-genus49-jacobian-computable-isogeny}，
\[
J(X)\sim_{\QQ}J(H)\times J(Z)^2\times J(C_F)^3\times P^3.
\]
同源在有理一维同调上诱导直和分解
\[
H_1\!\bigl(J(X)(\CC),\QQ\bigr)
\cong
H_1\!\bigl(J(H)(\CC),\QQ\bigr)
\oplus
H_1\!\bigl(J(Z)(\CC),\QQ\bigr)^{\oplus 2}
\oplus
H_1\!\bigl(J(C_F)(\CC),\QQ\bigr)^{\oplus 3}
\oplus
H_1\!\bigl(P(\CC),\QQ\bigr)^{\oplus 3}.
\]
又因 \(H_1(J(C)(\CC),\ZZ)\) 与 \(H_1(C(\CC),\ZZ)\) 具有相同秩，故
\[
\cdim\!\bigl(\Lambda(X)\bigr)
=
\cdim\!\bigl(\Lambda(H)\bigr)
+2\,\cdim\!\bigl(\Lambda(Z)\bigr)
+3\,\cdim\!\bigl(\Lambda(C_F)\bigr)
+3\,\cdim\!\bigl(\Lambda(P)\bigr).
\]
代入上式即得
\[
\cdim\!\bigl(\Lambda(X)\bigr)=10+2\cdot 8+3\cdot 6+3\cdot 18=98.
\]
另一方面，\(g(X)=49\)，故
\[
\cdim\!\bigl(\Lambda(X)\bigr)=2g(X)=98.
\]
两边一致，命题得证。
\end{proof}

\begin{conclusion}[缺失商 Jacobian 的显式因子化]\label{con:xi-time-part9n-c2-c3-quotient-jacobian-complete-splitting}
沿用上述 \(S_4\)-Hurwitz 商塔记号，并记
\[
W_3:=X/C_3,\qquad
W_2:=X/C_2,\qquad
P:=\operatorname{Prym}(Y/Z).
\]
则有显式同源分解
\[
\boxed{\;
J(W_2)\sim_{\QQ}J(Z)\times J(C_F)^2\times P,
\;}
\]
\[
\boxed{\;
J(W_3)\sim_{\QQ}J(H)\times J(C_F)\times P.
\;}
\]
特别地，
\[
g(W_2)=4+2\cdot 3+9=19,
\qquad
g(W_3)=5+3+9=17,
\]
与结论 \ref{con:xi-time-part9n-s4-all-subgroup-genera} 的亏格恢复完全一致。
\end{conclusion}
\begin{proof}
先证 \(C_2\)-商。由定理 \ref{thm:killo-s4-burnside-kani-rosen-prym-square}，
\[
J(X/C_2)\sim_{\QQ}J(C_F)^2\times J(Y).
\]
另一方面，\(Y\to Z\) 是二重覆盖，而 \(P=\operatorname{Prym}(Y/Z)\)。对二重覆盖的标准 Jacobian--Prym 分解给出
\[
J(Y)\sim_{\QQ}J(Z)\times P.
\]
代入即得
\[
J(X/C_2)\sim_{\QQ}J(Z)\times J(C_F)^2\times P.
\]

再证 \(C_3\)-商。由定理 \ref{thm:killo-s4-genus49-jacobian-computable-isogeny}，
\[
J(X)\sim_{\QQ}J(H)\times J(Z)^2\times J(C_F)^3\times P^3,
\]
其中四个可见因子分别对应于 \(\mathrm{sgn},V_2,V_3,V_3'\) 四个非同构有理不可约表示块。对任意子群 \(K\le S_4\)，商曲线的 holomorphic differential 空间满足
\[
H^0(X/K,\Omega_{X/K}^1)\cong H^0(X,\Omega_X^1)^K,
\]
故每个等型块在 \(J(X/K)\) 中出现的次数正是相应不可约表示的 \(K\)-不变维数。

对 \(K=C_3\)，由
\[
\rho_{C_3}=\mathbf 1\oplus \mathrm{sgn}\oplus V_3\oplus V_3'
\]
可知
\[
\bigl(\dim \mathrm{sgn}^{C_3},\dim V_2^{C_3},\dim V_3^{C_3},\dim (V_3')^{C_3}\bigr)=(1,0,1,1).
\]
因此 \(J(X/C_3)\) 恰由一个 \(\mathrm{sgn}\)-块、一个 \(V_3\)-块与一个 \(V_3'\)-块组成，从而
\[
J(X/C_3)\sim_{\QQ}J(H)\times J(C_F)\times P.
\]
最后，两式的维数分别为 \(4+2\cdot 3+9=19\) 与 \(5+3+9=17\)，恰与结论 \ref{con:xi-time-part9n-s4-all-subgroup-genera} 一致。证毕。
\end{proof}

\begin{conclusion}[良素处换位商与三循环商局部 \texorpdfstring{$L$}{L} 因子的完全分裂]\label{con:xi-time-part9n-c2-c3-local-lfactor-complete-splitting}
设 \(v\) 为一处使
\[
H,\ Z,\ C_F,\ Y,\ X/C_2,\ X/C_3
\]
及 Prym 因子 \(P=\operatorname{Prym}(Y/Z)\) 同时良约化的有限素位。记各自的局部 \(L\)-多项式分子为
\[
P_H(v,T),\quad P_Z(v,T),\quad P_{C_F}(v,T),\quad
P_P(v,T),\quad P_{X/C_2}(v,T),\quad P_{X/C_3}(v,T).
\]
则有完全乘积分解
\[
\boxed{\;
P_{X/C_2}(v,T)=P_Z(v,T)\,P_{C_F}(v,T)^2\,P_P(v,T),
\;}
\]
\[
\boxed{\;
P_{X/C_3}(v,T)=P_H(v,T)\,P_{C_F}(v,T)\,P_P(v,T).
\;}
\]
\end{conclusion}
\begin{proof}
由结论 \ref{con:xi-time-part9n-c2-c3-quotient-jacobian-complete-splitting}，
\[
J(X/C_2)\sim_{\QQ}J(Z)\times J(C_F)^2\times P,
\]
\[
J(X/C_3)\sim_{\QQ}J(H)\times J(C_F)\times P.
\]
对共同良约化素位 \(v\)，同源阿贝尔簇的 \(\ell\)-进 Tate 模半单化具有相同的 Frobenius 特征多项式，而直积对应特征多项式相乘。故两式立即推出
\[
P_{X/C_2}(v,T)=P_Z(v,T)\,P_{C_F}(v,T)^2\,P_P(v,T),
\]
\[
P_{X/C_3}(v,T)=P_H(v,T)\,P_{C_F}(v,T)\,P_P(v,T).
\]
证毕。
\end{proof}

\begin{conclusion}[两类 Prym 极化缺陷诱导的最小同源度刚性]\label{con:xi-time-part9n-prym-polarization-minimal-isogeny-degree}
记
\[
A:=\operatorname{Prym}(C_6/H),
\qquad
P:=\operatorname{Prym}(Y/Z),
\]
并令 \(\lambda_A,\lambda_P\) 为命题 \ref{prop:killo-two-prym-polarization-types} 给出的诱导极化。于是
\[
\mathrm{type}(\lambda_A)=(1,1,1,1,3,3,3,3),
\qquad
\mathrm{type}(\lambda_P)=(1,1,1,1,2,2,2,2,2).
\]
则成立：
\begin{enumerate}
  \item 若 \((B,\lambda_B)\) 为八维主极化阿贝尔簇，且存在同源
  \[
  f:B\to A
  \]
  与整数 \(n\ge 1\) 使
  \[
  f^*\lambda_A=n\lambda_B,
  \]
  则
  \[
  \deg f\ge 3^4=81,
  \]
  且等号当且仅当 \(n=3\)。
  \item 若 \((C,\lambda_C)\) 为九维主极化阿贝尔簇，且存在同源
  \[
  g:C\to P
  \]
  与整数 \(m\ge 1\) 使
  \[
  g^*\lambda_P=m\lambda_C,
  \]
  则
  \[
  \deg g\ge 2^4=16,
  \]
  且等号当且仅当 \(m=2\)。
\end{enumerate}
\end{conclusion}
\begin{proof}
由极化型公式
\[
\deg(\lambda)=\prod_{i=1}^{g}d_i^2,
\]
得
\[
\deg(\lambda_A)=3^8,
\qquad
\deg(\lambda_P)=2^{10}.
\]

先看 \(A\) 的情形。由于 \(\dim B=8\) 且 \(\lambda_B\) 为主极化，
\[
\deg(n\lambda_B)=n^{16}.
\]
另一方面，极化对同源的拉回满足
\[
\deg(f^*\lambda_A)=(\deg f)^2\deg(\lambda_A).
\]
由 \(f^*\lambda_A=n\lambda_B\) 得
\[
(\deg f)^2\cdot 3^8=n^{16},
\qquad\text{即}\qquad
(\deg f)^2=\frac{n^{16}}{3^8}.
\]
右端为整数平方，故其 \(3\)-进赋值必须为非负偶数。设 \(v_3(n)=a\)，则
\[
v_3\!\bigl((\deg f)^2\bigr)=16a-8.
\]
因此必有 \(a\ge 1\)，即 \(3\mid n\)。写 \(n=3r\)，则
\[
(\deg f)^2=3^8r^{16},
\qquad
\deg f=3^4r^8\ge 3^4.
\]
等号成立当且仅当 \(r=1\)，亦即 \(n=3\)。

再看 \(P\) 的情形。由于 \(\dim C=9\) 且 \(\lambda_C\) 为主极化，
\[
\deg(m\lambda_C)=m^{18}.
\]
同理，
\[
(\deg g)^2\cdot 2^{10}=m^{18},
\qquad\text{即}\qquad
(\deg g)^2=\frac{m^{18}}{2^{10}}.
\]
设 \(v_2(m)=b\)，则
\[
v_2\!\bigl((\deg g)^2\bigr)=18b-10.
\]
其为非负偶整数，故必有 \(b\ge 1\)，即 \(2\mid m\)。写 \(m=2s\)，则
\[
(\deg g)^2=2^8s^{18},
\qquad
\deg g=2^4s^9\ge 2^4.
\]
等号成立当且仅当 \(s=1\)，亦即 \(m=2\)。证毕。
\end{proof}

在 \(P\) 本身的极化缺陷之外，\(V_3'\)-等型块 \(A_{3'}\sim_{\QQ} P^3\) 的主极化化代价、全部额外有理端同态以及正维子块类型也可由同一分解立即锁定。

\begin{theorem}[Prym 三重等型块的主极化化、条件性矩阵端代数与子块分类]\label{thm:xi-time-part9n-prym-cubed-principalization-endomorphism}
\leanverified{paper\_xi\_time\_part9n\_prym\_cubed\_principalization\_endomorphism}
沿用结论 \ref{con:xi-time-part9n-prym-polarization-minimal-isogeny-degree} 的记号，记
\[
P:=\operatorname{Prym}(Y/Z),
\qquad
A_{3'}\sim_{\QQ} P^3,
\qquad
\mathrm{type}(\lambda_P)=(1^4,2^5).
\]
令 \(\lambda_{P^3}\) 表示 \(P^3\) 上的乘积极化。则成立：
\begin{enumerate}
  \item \(\lambda_{P^3}\) 的极化型为
  \[
  \mathrm{type}(\lambda_{P^3})=(1^{12},2^{15}).
  \]
  \item 若 \((B,\Theta)\) 为主极化阿贝尔簇，且
  \[
  f:P^3\to B
  \]
  为同源并满足
  \[
  f^*\Theta=\lambda_{P^3},
  \]
  则
  \[
  \deg f=2^{15}.
  \]
  特别地，\(P^3\) 的主极化化最小次数恰为 \(2^{15}\)。
  \item 若猜想 \ref{conj:killo-minimal-extra-endomorphism-large-prym} 的第一条成立，即 \(P\) 在 \(\overline{\QQ}\) 上绝对单，且
  \[
  \operatorname{End}_{\overline{\QQ}}(P)=\ZZ,
  \]
  则
  \[
  \operatorname{End}^0_{\QQ}(A_{3'})\cong M_3(\QQ),
  \qquad
  Z\!\bigl(\operatorname{End}^0_{\QQ}(A_{3'})\bigr)=\QQ.
  \]
  并且 \(A_{3'}\) 的任意正维 \(\QQ\)-阿贝尔子簇都同源于
  \[
  P^r,\qquad 1\le r\le 3.
  \]
  特别地，\(A_{3'}\) 的任意 \(9\) 维简单 \(\QQ\)-因子都唯一同源于 \(P\)。
\end{enumerate}
\end{theorem}
\begin{proof}
乘积极化的 elementary divisors 由各因子的 elementary divisors 直接并列给出，故
\[
\mathrm{type}(\lambda_{P^3})
=
(1^4,2^5)^{\times 3}
=(1^{12},2^{15}).
\]
这证明了第 \(1\) 条。

由极化型公式，
\[
\deg(\lambda_{P^3})=(2^{15})^2.
\]
另一方面，若 \(\Theta\) 为主极化且 \(f^*\Theta=\lambda_{P^3}\)，则
\[
\deg(\lambda_{P^3})
=
\deg(f^*\Theta)
=
(\deg f)^2\deg(\Theta)
=
(\deg f)^2.
\]
故
\[
\deg f=2^{15}.
\]
再者，\(\ker(\lambda_{P^3})\) 作为有限辛群含有阶为 \(2^{15}\) 的极大各向同性子群 \(K\)。对商映射
\[
q:P^3\to P^3/K
\]
应用极化下降的标准构造，可得某个主极化 \(\Theta_K\) 满足
\[
q^*\Theta_K=\lambda_{P^3}.
\]
因此 \(P^3\) 的主极化化确可由次数 \(2^{15}\) 的同源实现，从而该次数既是必要下界，也是可达下界，第 \(2\) 条成立。

若猜想 \ref{conj:killo-minimal-extra-endomorphism-large-prym} 的第一条成立，则
\[
\operatorname{End}^0_{\overline{\QQ}}(P)
=
\operatorname{End}_{\overline{\QQ}}(P)\otimes_{\ZZ}\QQ
=
\QQ.
\]
因而
\[
\operatorname{End}^0_{\QQ}(P)=\QQ.
\]
另一方面，由结论 \ref{con:xi-terminal-zm-s4-jacobian-a3p-isog-prym3-cubed}，
\[
A_{3'}\sim_{\QQ} P^3.
\]
同源阿贝尔簇的有理端同态代数在 \(\QQ\)-同源范畴中同构，而任意阿贝尔簇 \(A\) 都满足
\[
\operatorname{End}^0_{\QQ}(A^3)\cong
M_3\!\bigl(\operatorname{End}^0_{\QQ}(A)\bigr).
\]
取 \(A=P\) 即得
\[
\operatorname{End}^0_{\QQ}(A_{3'})
\cong
\operatorname{End}^0_{\QQ}(P^3)
\cong
M_3(\QQ).
\]
其中心显然等于 \(\QQ\)。

再设 \(C\subseteq A_{3'}\) 为任意正维 \(\QQ\)-阿贝尔子簇。经由 \(A_{3'}\sim_{\QQ} P^3\) 的同源识别，\(C\) 在 \(\QQ\)-同源范畴中对应于某个幂等元
\[
e\in \operatorname{End}^0_{\QQ}(P^3)\cong M_3(\QQ)
\]
的像。线性代数表明，\(M_3(\QQ)\) 中任一幂等元都按共轭同价于
\[
\operatorname{diag}(I_r,0)\qquad (0\le r\le 3),
\]
其中 \(r=\operatorname{rank}(e)\) 唯一确定。于是
\[
\operatorname{Im}(e)\sim_{\QQ} P^r.
\]
由于 \(C\) 正维，故 \(1\le r\le 3\)。这就给出了全部正维 \(\QQ\)-阿贝尔子簇的同源类型分类。

最后，若 \(C\) 还是简单且 \(\dim C=9=\dim P\)，则只能有 \(r=1\)，从而
\[
C\sim_{\QQ} P.
\]
故第 \(3\) 条全部成立。证毕。
\end{proof}

\endinput

在 \(V_3'\)-等型三重块 \(P^3\) 的主极化化最小次数与矩阵端代数已经锁定之后，还可把这一个 \(27\)-维可见块在 \(\QQ\)-同源范畴中的子等型与 Rosati 对称端代数完全显式化。

\begin{theorem}[\texorpdfstring{$27$}{27} 维 Prym 等型块的 \texorpdfstring{$\QQ$}{Q}-子等型完全分类]\label{thm:xi-time-part9n1e-prym-cubed-q-subisogeny-rosati-classification}
沿用
\[
J(X)\sim_{\QQ}J(H)\times J(Z)^2\times J(C_F)^3\times P^3,
\qquad
P:=\operatorname{Prym}(Y/Z),
\qquad
\dim P=9
\]
的记号，并进一步假设猜想 \ref{conj:killo-minimal-extra-endomorphism-large-prym} 的第一条成立，即
\[
P\ \text{在 }\overline{\QQ}\text{ 上绝对单，且}\ \operatorname{End}_{\overline{\QQ}}(P)=\ZZ.
\]
则成立：
\begin{enumerate}
  \item \textbf{（有理端代数）}
        \[
        \operatorname{End}_{\QQ}^0(P^3)\cong M_3(\QQ).
        \]
  \item \textbf{（\(\QQ\)-子等型完全分类）}
        \(P^3\) 的一切 \(\QQ\)-阿贝尔子簇，在 \(\QQ\)-同源意义下恰为
        \[
        P^r\qquad (r=0,1,2,3)
        \]
        这四种类型；不存在任何额外的、非幂次型的 \(\QQ\)-子等型。
  \item \textbf{（Rosati 对称端代数）}
        若在 \(P^3\) 上取乘积极化 \(\lambda_{P^3}\)，并记相应 Rosati 反自同态为 \(\dagger\)，则
        \[
        \operatorname{End}_{\QQ}^0(P^3)^{\dagger}
        \cong
        \operatorname{Sym}_3(\QQ).
        \]
        因而其有理极化锥恰由正定对称 \(3\times 3\) 有理矩阵参数化。
\end{enumerate}
\end{theorem}
\begin{proof}
由假设
\[
\operatorname{End}_{\overline{\QQ}}(P)=\ZZ
\]
可知
\[
\operatorname{End}_{\overline{\QQ}}^0(P)=\QQ.
\]
而 \(\operatorname{End}_{\QQ}(P)\) 自然嵌入 \(\operatorname{End}_{\overline{\QQ}}(P)\)，故
\[
\QQ\subseteq \operatorname{End}_{\QQ}^0(P)\subseteq \operatorname{End}_{\overline{\QQ}}^0(P)=\QQ,
\]
从而
\[
\operatorname{End}_{\QQ}^0(P)=\QQ.
\]
对任意阿贝尔簇 \(A\) 与整数 \(m\ge 1\)，标准等型公式给出
\[
\operatorname{End}_{\QQ}^0(A^m)\cong M_m\!\bigl(\operatorname{End}_{\QQ}^0(A)\bigr).
\]
取 \(A=P\)、\(m=3\) 即得
\[
\operatorname{End}_{\QQ}^0(P^3)\cong M_3(\QQ),
\]
证明了第 \(1\) 条。

现证第 \(2\) 条。阿贝尔簇的 \(\QQ\)-同源范畴是半单伪阿贝尔范畴，因此 \(P^3\) 的 \(\QQ\)-阿贝尔子簇的同源类与 \(\operatorname{End}_{\QQ}^0(P^3)\) 中的幂等元一一对应。借助第 \(1\) 条，可把此问题转化为 \(M_3(\QQ)\) 中幂等矩阵的分类。

对任意 \(e\in M_3(\QQ)\) 满足 \(e^2=e\)，其最小多项式整除 \(t(t-1)\)，故无重根，从而 \(e\) 在 \(\QQ\) 上可对角化。于是 \(e\) 必与某个
\[
\operatorname{diag}(I_r,0)\qquad (r=0,1,2,3)
\]
在 \(\operatorname{GL}_3(\QQ)\) 下共轭。对应的像对象在 \(\QQ\)-同源意义下正是 \(P^r\)。反过来，这四个秩值显然都可实现。故 \(P^3\) 的 \(\QQ\)-子等型只有
\[
0,\ P,\ P^2,\ P^3
\]
四类，而不存在任何额外的非幂次型子等型。

最后处理第 \(3\) 条。由乘积极化 \(\lambda_{P^3}=\lambda_P^{\times 3}\) 所诱导的 Rosati 反自同态，在
\[
\operatorname{End}_{\QQ}^0(P^3)\cong M_3\!\bigl(\operatorname{End}_{\QQ}^0(P)\bigr)=M_3(\QQ)
\]
上的作用与“系数域上的 Rosati 对合加矩阵转置”一致。由于 \(\operatorname{End}_{\QQ}^0(P)=\QQ\) 且 \(\QQ\) 上唯一正反演为恒等，故 \(\dagger\) 在 \(M_3(\QQ)\) 上即与转置等价。因此
\[
\operatorname{End}_{\QQ}^0(P^3)^{\dagger}
\cong
\{M\in M_3(\QQ):M^{\mathsf T}=M\}
=
\operatorname{Sym}_3(\QQ).
\]
在此识别下，Rosati 正元正对应于正定对称矩阵，于是有理极化锥恰由正定对称 \(3\times 3\) 有理矩阵参数化。证毕。
\end{proof}

\noindent
因此，在大维 Prym 因子的绝对单纯性与端同态极小化假设下，\(P^3\) 不仅具有矩阵型有理端代数，而且其 \(\QQ\)-子等型、Rosati 对称端代数与极化锥都被 \(M_3(\QQ)\) 的线性代数完全刚性化。

\endinput

\paragraph{规范化纤维积、Hodge 张量分层与局部 Euler 因子的平方/立方刚性}
在 \(S_4\)-Hurwitz 商塔的亏格、总分歧、Jacobian 同源分解与 Prym 极化缺陷已经完全确定之后，还可继续把 \(A_4\)、\(D_8\)、\(V_4\) 三层商曲线之间的几何拼接、全纯微分的表示论来源以及良素处局部 Euler 因子的幂次结构进一步压缩为下列四条后果。

\begin{conclusion}[\texorpdfstring{$V_4$}{V4} 商曲线的规范化纤维积刻画]\label{con:xi-time-part9n-s4-c6-normalized-fiberproduct}
沿用 \(S_4\)-Hurwitz 商塔记号，并令
\[
H\longrightarrow \PP^1_u,\qquad Z\longrightarrow \PP^1_u
\]
为由 \(X\to \PP^1_u\) 诱导的自然映射。则有规范化纤维积识别
\[
\boxed{\;
C_6\ \cong\ \widetilde{H\times_{\PP^1_u} Z},
\;}
\]
其中右侧表示纤维积的规范化。

并且两条自然投影满足
\[
\deg\!\bigl(C_6\to H\bigr)=3,\qquad \deg\!\bigl(C_6\to Z\bigr)=2.
\]
其中 \(C_6\to H\) 为无分歧循环三重覆盖，而 \(C_6\to Z\) 的总分歧度恰为
\[
\boxed{\;12.\;}
\]
\end{conclusion}
\begin{proof}
在 \(S_4\) 中有群论恒等式
\[
A_4\cap D_8=V_4,\qquad A_4D_8=S_4.
\]
因此对正则 \(S_4\)-覆盖 \(X\to \PP^1_u\)，其函数域满足
\[
\QQ(H)=\QQ(X)^{A_4},\qquad
\QQ(Z)=\QQ(X)^{D_8},
\]
而在 \(\QQ(X)\) 中的复合域为
\[
\QQ(H)\,\QQ(Z)
=
\QQ(X)^{A_4\cap D_8}
=
\QQ(X)^{V_4}
=
\QQ(C_6).
\]
另一方面，\(\widetilde{H\times_{\PP^1_u} Z}\) 的函数域正是 \(\QQ(H)\) 与 \(\QQ(Z)\) 在 \(\QQ(X)\) 中的复合域，故
\[
\widetilde{H\times_{\PP^1_u} Z}\cong C_6.
\]

两条投影的次数由子群指数直接给出：
\[
[A_4:V_4]=3,\qquad [D_8:V_4]=2.
\]
又由定理 \ref{thm:killo-s4-even-subgroup-free-action-unramified-a4-lift}，
\[
C_6\to H
\]
本身就是无分歧循环三重覆盖。

对另一条二重映射，由结论 \ref{con:xi-time-part9n-s4-all-subgroup-genera}，
\[
g(C_6)=13,\qquad g(Z)=4.
\]
代入 Riemann--Hurwitz 公式
\[
2g(C_6)-2=2\bigl(2g(Z)-2\bigr)+R
\]
即得
\[
24=12+R,
\]
故 \(R=12\)。证毕。
\end{proof}

\begin{conclusion}[全纯微分空间的 \texorpdfstring{$S_4$}{S4} 张量型 Hodge 分层]\label{con:xi-time-part9n-s4-hodge-tensor-factorization}
记
\[
V:=H^0(X,\Omega_X^1).
\]
则 \(V\) 作为 \(S_4\)-表示与 Hodge 余切空间同时满足严格张量分解
\[
\boxed{\;
H^0(X,\Omega_X^1)
\cong
\mathrm{sgn}\otimes H^0(H,\Omega_H^1)
\oplus
V_2\otimes H^0(Z,\Omega_Z^1)
\oplus
V_3\otimes H^0(C_F,\Omega_{C_F}^1)
\oplus
V_3'\otimes H^0(P,\Omega_P^1).
\;}
\]
特别地，\(V_3'\)-等型部分全部由九维 Prym 因子 \(P\) 承载，不再混入任何来自曲线商 \(H\)、\(Z\)、\(C_F\) 的贡献。
\end{conclusion}
\begin{proof}
由定理 \ref{thm:killo-s4-genus49-jacobian-computable-isogeny}，
\[
J(X)\sim_{\QQ}J(H)\times J(Z)^2\times J(C_F)^3\times P^3.
\]
对任意阿贝尔簇同源分解取单位元处余切空间，并用 Jacobian 的标准识别
\[
H^0\!\bigl(J(C),\Omega_{J(C)}^1\bigr)\cong H^0(C,\Omega_C^1),
\]
得到向量空间分解
\[
H^0(X,\Omega_X^1)
\cong
H^0(H,\Omega_H^1)
\oplus
H^0(Z,\Omega_Z^1)^{\oplus 2}
\oplus
H^0(C_F,\Omega_{C_F}^1)^{\oplus 3}
\oplus
H^0(P,\Omega_P^1)^{\oplus 3}.
\]
其中
\[
\dim H^0(H,\Omega_H^1)=5,\qquad
\dim H^0(Z,\Omega_Z^1)=4,
\]
\[
\dim H^0(C_F,\Omega_{C_F}^1)=3,\qquad
\dim H^0(P,\Omega_P^1)=9.
\]

另一方面，结论 \ref{con:xi-time-part9n-s4-all-subgroup-genera} 给出的 Chevalley--Weil 分解为
\[
H^0(X,\Omega_X^1)
\cong
5\cdot \mathrm{sgn}\oplus 4\cdot V_2\oplus 3\cdot V_3\oplus 9\cdot V_3',
\]
故四个等型块的总维数分别为
\[
1\cdot 5,\qquad 2\cdot 4,\qquad 3\cdot 3,\qquad 3\cdot 9.
\]
它们恰与
\[
J(H),\qquad J(Z)^2,\qquad J(C_F)^3,\qquad P^3
\]
在余切空间中的维数分块完全一致。由于 \(\mathrm{sgn}\)、\(V_2\)、\(V_3\)、\(V_3'\) 两两不同构，等型分解的重数空间唯一，于是只能分别识别为
\[
H^0(H,\Omega_H^1),\qquad
H^0(Z,\Omega_Z^1),\qquad
H^0(C_F,\Omega_{C_F}^1),\qquad
H^0(P,\Omega_P^1).
\]
这就得到所示张量型分解。最后一条断言只是把 \(V_3'\)-块与 \(P^3\) 的唯一对应关系重新表述。证毕。
\end{proof}

\begin{conclusion}[良素处局部 Hasse--Weil 分子的平方/立方块与点数整除律]\label{con:xi-time-part9n-s4-local-euler-square-cube}
设 \(p\) 为一个使
\[
X,\ H,\ C_6,\ Z,\ C_F
\]
及阿贝尔簇 \(P\) 同时良约化的素数。记各自的局部 Hasse--Weil 分子为
\[
P_{H,p}(T),\quad P_{C_6,p}(T),\quad P_{Z,p}(T),\quad
P_{C_F,p}(T),\quad P_{X,p}(T),\quad P_{P,p}(T).
\]
则有严格因子分解
\[
\boxed{\;
P_{C_6,p}(T)=P_{H,p}(T)\,P_{Z,p}(T)^2,
\;}
\]
\[
\boxed{\;
P_{X,p}(T)=P_{H,p}(T)\,P_{Z,p}(T)^2\,P_{C_F,p}(T)^3\,P_{P,p}(T)^3.
\;}
\]
等价地，
\[
P_{X,p}(T)=P_{H,p}(T)\,P_{Z,p}(T)^2
\bigl(P_{C_F,p}(T)P_{P,p}(T)\bigr)^3,
\]
因此 \(C_6\) 的局部分子必带严格平方块，而 \(X\) 的局部分子必带严格立方块。

若记
\[
t_p(P):=
\operatorname{Tr}\!\Bigl(
\mathrm{Frob}_p\mid H^1_{\text{ét}}(P_{\overline{\FF}_p},\QQ_\ell)
\Bigr),
\]
则进一步有点数公式
\[
\boxed{\;
\#C_6(\FF_p)=\#H(\FF_p)+2\,\#Z(\FF_p)-2(p+1),
\;}
\]
\[
\boxed{\;
\#X(\FF_p)=\#H(\FF_p)+2\,\#Z(\FF_p)+3\,\#C_F(\FF_p)-5(p+1)-3\,t_p(P).
\;}
\]
特别地，
\[
\boxed{\;
\#X(\FF_p)-\#H(\FF_p)-2\,\#Z(\FF_p)-3\,\#C_F(\FF_p)+5(p+1)\equiv 0\pmod 3.
\;}
\]
\end{conclusion}
\begin{proof}
由定理 \ref{thm:killo-s4-burnside-kani-rosen-prym-square}，
\[
J(C_6)\sim_{\QQ}J(H)\times J(Z)^2,
\]
而由定理 \ref{thm:killo-s4-genus49-jacobian-computable-isogeny}，
\[
J(X)\sim_{\QQ}J(H)\times J(Z)^2\times J(C_F)^3\times P^3.
\]
在共同良约化素数 \(p\) 处，同源阿贝尔簇的 \(\ell\)-进 \(H^1\) 具有相同 Frobenius 特征多项式，而直积对应特征多项式相乘，故立即得到两条局部分子分解。

对任意曲线 \(C/\FF_p\)，写
\[
P_{C,p}(T)=1-a_p(C)T+\cdots,\qquad a_p(C)=p+1-\#C(\FF_p).
\]
对阿贝尔簇 \(P\)，同理有
\[
P_{P,p}(T)=1-t_p(P)T+\cdots.
\]
比较
\[
P_{C_6,p}(T)=P_{H,p}(T)\,P_{Z,p}(T)^2
\]
的一次项，得到
\[
a_p(C_6)=a_p(H)+2a_p(Z),
\]
即
\[
\#C_6(\FF_p)=\#H(\FF_p)+2\,\#Z(\FF_p)-2(p+1).
\]

再比较
\[
P_{X,p}(T)=P_{H,p}(T)\,P_{Z,p}(T)^2\,P_{C_F,p}(T)^3\,P_{P,p}(T)^3
\]
的一次项，得到
\[
a_p(X)=a_p(H)+2a_p(Z)+3a_p(C_F)+3t_p(P).
\]
把 \(a_p(C)=p+1-\#C(\FF_p)\) 代回并整理，即得
\[
\#X(\FF_p)=\#H(\FF_p)+2\,\#Z(\FF_p)+3\,\#C_F(\FF_p)-5(p+1)-3\,t_p(P).
\]
最后一条同余式只是把上式移项后按模 \(3\) 约化。证毕。
\end{proof}

\begin{conclusion}[双平方块上的 \texorpdfstring{$\QQ(\zeta_3)$}{Q(zeta3)}-矩阵乘法与 cyclotomic norm-square]\label{con:xi-time-part9n-s4-jz2-cyclotomic-norm-square}
记
\[
K:=\QQ(\zeta_3),\qquad A:=\operatorname{Prym}(C_6/H).
\]
则存在 \(\QQ\)-代数嵌入
\[
\boxed{\;
K\hookrightarrow \operatorname{End}^0\!\bigl(J(Z)^2\bigr)
\cong
M_2\!\bigl(\operatorname{End}^0(J(Z))\bigr).
\;}
\]
因此亏格 \(4\) 因子 \(J(Z)\) 的平方块带有必然的 cyclotomic matrix multiplication。

进一步地，对任意 \(Z\) 的良素数 \(p\)，存在次数 \(8\) 的多项式
\[
Q_{Z,p}(T)\in K[T]
\]
使得
\[
\boxed{\;
P_{Z,p}(T)^2=
\operatorname{Norm}_{K/\QQ}\!\bigl(Q_{Z,p}(T)\bigr).
\;}
\]
换言之，\(P_{Z,p}(T)^2\) 不是任意平方，而是一个 cyclotomic norm-square。
\end{conclusion}
\begin{proof}
由命题 \ref{prop:killo-two-prym-polarization-types}，
\[
K=\QQ(\zeta_3)\subset \operatorname{End}^0(A).
\]
再由定理 \ref{thm:killo-s4-burnside-kani-rosen-prym-square}，
\[
A\sim_{\QQ}J(Z)^2.
\]
沿任一 \(\QQ\)-同源把端同态代数共轭传送，便得到
\[
K\hookrightarrow \operatorname{End}^0\!\bigl(J(Z)^2\bigr).
\]
而
\[
\operatorname{End}^0\!\bigl(J(Z)^2\bigr)\cong M_2\!\bigl(\operatorname{End}^0(J(Z))\bigr)
\]
是标准矩阵代数识别，故第一条结论成立。

现取 \(Z\) 的一处良素数 \(p\) 与 \(\ell\neq p\)。由于 \(K\) 的作用来自代数端同态，它与 Frobenius 作用对易，因此
\[
V_\ell\!\bigl(J(Z)^2\bigr):=
H^1_{\text{ét}}\!\bigl(J(Z)^2_{\overline{\FF}_p},\QQ_\ell\bigr)
\]
可视为一个 \(K\otimes_{\QQ}\QQ_\ell\)-线性表示。其 \(\QQ_\ell\)-维数为 \(16\)，故作为 \(K\)-线性空间的维数为 \(8\)。记
\[
Q_{Z,p}(T)\in K[T]
\]
为 Frobenius 在该 \(K\)-线性表示上的特征多项式，则 \(\deg Q_{Z,p}=8\)。

按定义，\(\operatorname{Norm}_{K/\QQ}(Q_{Z,p}(T))\) 正是 Frobenius 在同一空间上作为 \(\QQ_\ell\)-线性变换的特征多项式。另一方面，
\[
H^1_{\text{ét}}\!\bigl(J(Z)^2_{\overline{\FF}_p},\QQ_\ell\bigr)
\cong
H^1_{\text{ét}}\!\bigl(Z_{\overline{\FF}_p},\QQ_\ell\bigr)^{\oplus 2},
\]
故该 \(\QQ_\ell\)-线性特征多项式恰为 \(P_{Z,p}(T)^2\)。于是
\[
P_{Z,p}(T)^2=
\operatorname{Norm}_{K/\QQ}\!\bigl(Q_{Z,p}(T)\bigr).
\]
证毕。
\end{proof}

\begin{conclusion}[Hurwitz 塔中的素数分离型 Prym 刚性]\label{con:xi-time-part9n-prym-prime-separated-rigidity}
记
\[
A_{(3)}:=\operatorname{Prym}(C_6/H),
\qquad
A_{(2)}:=P=\operatorname{Prym}(Y/Z).
\]
并记 \(\lambda_{A_{(3)}}\)、\(\lambda_{A_{(2)}}\) 为其自然诱导极化。则三次层与二次层的 Prym 极化缺陷严格按素数分离：
\[
\boxed{\;
A_{(3)}\sim_{\QQ}J(Z)^2,\qquad
\mathrm{type}(\lambda_{A_{(3)}})=(1,1,1,1,3,3,3,3),\qquad
\QQ(\zeta_3)\subset \operatorname{End}^0(A_{(3)}).\;}
\]
\[
\boxed{\;
\mathrm{type}(\lambda_{A_{(2)}})=(1,1,1,1,2,2,2,2,2).\;}
\]
特别地，三次分支所产生的全部非主极化缺陷都局限于 \(3\)-初等方向，而二次分支所产生的全部非主极化缺陷都局限于 \(2\)-初等方向；两者在极化型与端同态约束上互不混杂。
\end{conclusion}
\begin{proof}
由定理 \ref{thm:killo-s4-burnside-kani-rosen-prym-square}，
\[
\operatorname{Prym}(C_6/H)\sim_{\QQ}J(Z)^2.
\]
再由命题 \ref{prop:killo-two-prym-polarization-types}，
\[
\mathrm{type}(\lambda_{A_{(3)}})=(1,1,1,1,3,3,3,3),
\qquad
\mathrm{type}(\lambda_{A_{(2)}})=(1,1,1,1,2,2,2,2,2),
\]
并且
\[
\QQ(\zeta_3)\subset \operatorname{End}^0\!\bigl(\operatorname{Prym}(C_6/H)\bigr).
\]
这与结论 \ref{con:xi-time-part9n-s4-jz2-cyclotomic-norm-square} 中由 \(J(Z)^2\) 读出的 cyclotomic matrix multiplication 完全一致。故 \(A_{(3)}\) 的全部非单位初等因子都是 \(3\)，而 \(A_{(2)}\) 的全部非单位初等因子都是 \(2\)，从而两类 Prym 的非主极化缺陷严格分裂为互素的两块。证毕。
\end{proof}

\begin{theorem}[\texorpdfstring{$\operatorname{Prym}(C_6/H)\to J(Z)^2$}{Prym(C6/H) to J(Z)^2} 的极化兼容同源度刚性]\label{thm:xi-time-part9n1a-prym-c6h-jz2-polarized-isogeny-degree}
\leanverified{paper\_xi\_time\_part9n1a\_prym\_c6h\_jz2\_polarized\_isogeny\_degree}
记
\[
A:=\operatorname{Prym}(C_6/H),
\qquad
B:=J(Z)^2.
\]
令 \(\lambda_A\) 为 \(A\) 的自然诱导极化，\(\lambda_B\) 为 \(B\) 上的乘积主极化。若
\[
f:A\to B
\]
是满足
\[
f^*\lambda_B=\lambda_A
\]
的同源，则
\[
\deg f=3^4.
\]
\end{theorem}
\begin{proof}
由定理 \ref{thm:killo-s4-burnside-kani-rosen-prym-square}，
\[
A\sim_{\QQ}B.
\]
再由命题 \ref{prop:killo-two-prym-polarization-types}，
\[
\mathrm{type}(\lambda_A)=(1,1,1,1,3,3,3,3).
\]
因此
\[
\deg(\lambda_A)
=
(1\cdot1\cdot1\cdot1\cdot3\cdot3\cdot3\cdot3)^2
=
3^8,
\qquad
\deg(\lambda_B)=1.
\]
若 \(f^*\lambda_B=\lambda_A\)，则极化拉回的度数公式给出
\[
\deg(\lambda_A)
=
\deg(f^*\lambda_B)
=
(\deg f)^2\deg(\lambda_B)
=
(\deg f)^2.
\]
故
\[
(\deg f)^2=3^8,
\]
从而
\[
\deg f=3^4.
\]
证毕。
\end{proof}

\begin{corollary}[\texorpdfstring{$V_2$}{V2}-等型 Jacobian 块的 Eisenstein 乘法]\label{cor:xi-time-part9n1a-v2-block-eisenstein-multiplication}
记
\[
A_{V_2}:=J(Z)^2.
\]
则有
\[
\QQ(\zeta_3)\subset \operatorname{End}^0\!\bigl(A_{V_2}\bigr)
\subset
\operatorname{End}^0_{\QQ}\!\bigl(J(X)\bigr).
\]
特别地，属 \(49\) 的 \(S_4\)-Hurwitz Jacobian 分解中，对应 \(V_2\)-等型的八维块天然携带一条 Eisenstein 乘法结构，而不再是纯实型块。
\end{corollary}
\begin{proof}
第一重包含由结论 \ref{con:xi-time-part9n-s4-jz2-cyclotomic-norm-square} 直接给出。另一方面，定理 \ref{thm:xi-time-part9n1b-jacobian-projector-idempotent-algebra} 已把 \(J(Z)^2\) 识别为 \(J(X)\) 的一个规范块对角等型分量，故其有理端同态代数自然嵌入
\[
\operatorname{End}^0_{\QQ}\!\bigl(J(X)\bigr).
\]
由于 \(\QQ(\zeta_3)\) 是虚二次域，所述 \(V_2\)-块遂带有 Eisenstein 乘法。证毕。
\end{proof}

\endinput

\paragraph{Prym--Chevalley--Weil 互反、商曲线 genus 反演与投影代数刚性}
在 \(S_4\)-Hurwitz 商塔的亏格恢复、Hodge 张量分层、Burnside--Kani--Rosen 同源关系与 Jacobian 分块已经明确之后，还可继续把其中若干隐含的互反与刚性压缩为下列五条后果。

\begin{theorem}[\texorpdfstring{$S_4$}{S4}-几何表示的无余因子实现]\label{thm:xi-time-part9n1b-prym-chevalley-weil-reciprocity}
记
\[
V:=H^0(X,\Omega_X^1),
\qquad
A_{\mathrm{sgn}}:=J(H),\quad
A_{V_2}:=J(Z),\quad
A_{V_3}:=J(C_F),\quad
A_{V_3'}:=P.
\]
则四个非平凡有理不可约表示
\[
\mathrm{sgn},\qquad V_2,\qquad V_3,\qquad V_3'
\]
分别且唯一地由几何因子
\[
J(H),\qquad J(Z),\qquad J(C_F),\qquad P
\]
承载。更具体地，表示重数与几何维数满足精确互反
\[
\boxed{\;
V\cong
g(H)\cdot \mathrm{sgn}\ \oplus\
g(Z)\cdot V_2\ \oplus\
g(C_F)\cdot V_3\ \oplus\
\dim(P)\cdot V_3'.\;}
\]
并且 \(J(X)\) 的有理同源分解可重写为
\[
\boxed{\;
J(X)\sim_{\QQ}
A_{\mathrm{sgn}}^{\dim(\mathrm{sgn})}\times
A_{V_2}^{\dim(V_2)}\times
A_{V_3}^{\dim(V_3)}\times
A_{V_3'}^{\dim(V_3')}. \;}
\]
因此，\(S_4\)-微分表示的全部非平凡部分已被“商 Jacobian 因子 \(+\) 单个 Prym 因子”完全穷尽，不存在任何额外的隐藏等型几何块。
\end{theorem}
\begin{proof}
由结论 \ref{con:xi-time-part9n-s4-hodge-tensor-factorization}，
\[
V\cong
\mathrm{sgn}\otimes H^0(H,\Omega_H^1)
\oplus
V_2\otimes H^0(Z,\Omega_Z^1)
\oplus
V_3\otimes H^0(C_F,\Omega_{C_F}^1)
\oplus
V_3'\otimes H^0(P,\Omega_P^1).
\]
对四个重数空间分别取维数，并用
\[
\dim H^0(H,\Omega_H^1)=g(H),\quad
\dim H^0(Z,\Omega_Z^1)=g(Z),\quad
\dim H^0(C_F,\Omega_{C_F}^1)=g(C_F),\quad
\dim H^0(P,\Omega_P^1)=\dim(P),
\]
便得到第一式。

另一方面，由定理 \ref{thm:killo-s4-genus49-jacobian-computable-isogeny}，
\[
J(X)\sim_{\QQ}J(H)\times J(Z)^2\times J(C_F)^3\times P^3.
\]
再代入
\[
\dim(\mathrm{sgn})=1,\qquad \dim(V_2)=2,\qquad \dim(V_3)=\dim(V_3')=3,
\]
即得第二式。由于 \(S_4\) 的非平凡有理不可约表示恰止于这四类，而相应几何因子的维数分别为 \(5,4,3,9\)，总维数
\[
1\cdot 5+2\cdot 4+3\cdot 3+3\cdot 9=49=g(X)
\]
正好充满 \(J(X)\) 的全部 \(H^1\) 贡献，故不可能再有第五个隐藏等型块。证毕。
\end{proof}

\begin{corollary}[四个 genus 数据对非平凡 \(S_4\)-Hodge 向量的唯一恢复]\label{cor:xi-time-part9n1b-four-genera-recover-hodge-vector}
记
\[
V\cong
m_{\mathrm{sgn}}\cdot \mathrm{sgn}\oplus
m_{V_2}\cdot V_2\oplus
m_{V_3}\cdot V_3\oplus
m_{V_3'}\cdot V_3'.
\]
则仅由四个几何量
\[
g(X),\qquad g(H),\qquad g(C_6),\qquad g(C_F)
\]
便可唯一恢复重数向量 \((m_{\mathrm{sgn}},m_{V_2},m_{V_3},m_{V_3'})\)，其显式公式为
\[
\boxed{\;
m_{\mathrm{sgn}}=g(H),\qquad
m_{V_2}=\frac{g(C_6)-g(H)}{2},\qquad
m_{V_3}=g(C_F),\qquad
m_{V_3'}=\frac{g(X)-g(C_6)-3g(C_F)}{3}. \;}
\]
特别地，在本文数据
\[
g(X)=49,\qquad g(H)=5,\qquad g(C_6)=13,\qquad g(C_F)=3
\]
下，有
\[
(m_{\mathrm{sgn}},m_{V_2},m_{V_3},m_{V_3'})=(5,4,3,9).
\]
\end{corollary}
\begin{proof}
由定理 \ref{thm:killo-s4-burnside-kani-rosen-prym-square}，
\[
J(C_6)\sim_{\QQ}J(H)\times J(Z)^2,
\]
故取维数得到
\[
g(C_6)=g(H)+2g(Z),
\qquad
g(Z)=\frac{g(C_6)-g(H)}{2}.
\]
再由定理 \ref{thm:xi-time-part9n1b-prym-chevalley-weil-reciprocity}，
\[
m_{\mathrm{sgn}}=g(H),\qquad
m_{V_2}=g(Z),\qquad
m_{V_3}=g(C_F).
\]
最后仍由该定理，
\[
g(X)=\dim V=m_{\mathrm{sgn}}+2m_{V_2}+3m_{V_3}+3m_{V_3'},
\]
代入前三项即可解得
\[
m_{V_3'}=\frac{g(X)-g(C_6)-3g(C_F)}{3}.
\]
将 \((49,5,13,3)\) 代入即得 \((5,4,3,9)\)。证毕。
\end{proof}

\begin{theorem}[Burnside--Kani--Rosen 关系所迫定的两个隐含商曲线 genus]\label{thm:xi-time-part9n1b-kani-rosen-hidden-quotient-genera}
\leanverified{paper\_xi\_time\_part9n1b\_kani\_rosen\_hidden\_quotient\_genera}
沿用记号
\[
H:=X/A_4,\qquad C_6:=X/V_4,\qquad C_F:=X/S_3,\qquad Y:=X/C_4.
\]
若
\[
g(H)=5,\qquad g(C_6)=13,\qquad g(C_F)=3,\qquad g(Y)=13,
\]
则 Burnside--Kani--Rosen 关系在维数层面强制
\[
\boxed{\;
g(X/C_2)=19,\qquad g(X/C_3)=17.\;}
\]
\end{theorem}
\begin{proof}
由定理 \ref{thm:killo-s4-burnside-kani-rosen-prym-square}，
\[
J(X/C_2)\sim_{\QQ}J(C_F)^2\times J(Y).
\]
取维数即得
\[
g(X/C_2)=2g(C_F)+g(Y)=2\cdot 3+13=19.
\]

同一 Kani--Rosen 关系还给出
\[
J(C_6)\times J(X/C_3)^2\sim_{\QQ}J(H)^3\times J(C_F)^2\times J(Y)^2.
\]
再次取维数，得到
\[
g(C_6)+2g(X/C_3)=3g(H)+2g(C_F)+2g(Y).
\]
代入
\[
g(C_6)=13,\qquad g(H)=5,\qquad g(C_F)=3,\qquad g(Y)=13
\]
便有
\[
13+2g(X/C_3)=3\cdot 5+2\cdot 3+2\cdot 13=47,
\]
从而
\[
g(X/C_3)=17.
\]
证毕。
\end{proof}

\begin{theorem}[Jacobian 投影代数的显式幂等分层]\label{thm:xi-time-part9n1b-jacobian-projector-idempotent-algebra}
\leanverified{paper\_xi\_time\_part9n1b\_jacobian\_projector\_idempotent\_algebra}
有理自同态代数 \(\operatorname{End}^0_{\QQ}(J(X))\) 含有规范半单 \(\QQ\)-子代数
\[
\boxed{\;
\mathcal A_X:=\QQ\times M_2(\QQ)\times M_3(\QQ)\times M_3(\QQ)
\subseteq \operatorname{End}^0_{\QQ}(J(X)).\;}
\]
因而
\[
\boxed{\;\dim_{\QQ}\mathcal A_X=1+4+9+9=23.\;}
\]
此外，在 \(\mathcal A_X\) 中可取 \(9\) 个两两正交、非零且和为 \(1\) 的幂等元；并且存在 \(4\) 个中心幂等元，分别投影到
\[
J(H),\qquad J(Z)^2,\qquad J(C_F)^3,\qquad P^3
\]
这四个可见等型块。
\end{theorem}
\begin{proof}
定理 \ref{thm:xi-time-part9n-jacobian-endomorphism-lower-bound} 已给出块对角嵌入
\[
\QQ\times M_2(\QQ)\times M_3(\QQ)\times M_3(\QQ)
\hookrightarrow
\operatorname{End}^0_{\QQ}(J(X)),
\]
故 \(\mathcal A_X\) 的存在与 \(23\) 维结论立即成立。

现取标准矩阵单位
\[
E_{11}^{(2)},E_{22}^{(2)}\in M_2(\QQ),
\qquad
E_{11}^{(3)},E_{22}^{(3)},E_{33}^{(3)}\in M_3(\QQ),
\]
并在第二个 \(M_3(\QQ)\) 块上再取一份同样的对角矩阵单位。于是
\[
(1,0,0,0),
\quad
(0,E_{11}^{(2)},0,0),\ (0,E_{22}^{(2)},0,0),
\]
\[
(0,0,E_{11}^{(3)},0),\ (0,0,E_{22}^{(3)},0),\ (0,0,E_{33}^{(3)},0),
\]
\[
(0,0,0,E_{11}^{(3)}),\ (0,0,0,E_{22}^{(3)}),\ (0,0,0,E_{33}^{(3)})
\]
是 \(9\) 个两两正交的非零幂等元，并且其和恰为 \(\mathcal A_X\) 的单位元。

另一方面，四个块单位
\[
e_{\mathrm{sgn}}:=(1,0,0,0),\qquad
e_{V_2}:=(0,I_2,0,0),\qquad
e_{V_3}:=(0,0,I_3,0),\qquad
e_{V_3'}:=(0,0,0,I_3)
\]
位于直积代数的中心，并满足
\[
e_{\mathrm{sgn}}+e_{V_2}+e_{V_3}+e_{V_3'}=1.
\]
在定理 \ref{thm:killo-s4-genus49-jacobian-computable-isogeny} 给出的分块
\[
J(X)\sim_{\QQ}J(H)\times J(Z)^2\times J(C_F)^3\times P^3
\]
下，它们分别投影到所述四个可见等型块。证毕。
\end{proof}

\begin{corollary}[\texorpdfstring{$\operatorname{Prym}(C_6/H)$}{Prym(C6/H)} 在可见因子表中的唯一平方根]\label{cor:xi-time-part9n1b-prym-c6h-unique-square-root}
在几何因子表
\[
\{\,J(H),\ J(Z),\ J(C_F),\ P\,\}
\]
中，若某个因子 \(A\) 满足
\[
A^2\sim_{\QQ}\operatorname{Prym}(C_6/H),
\]
则必有
\[
\boxed{\;A\sim_{\QQ}J(Z).\;}
\]
\end{corollary}
\begin{proof}
由定理 \ref{thm:killo-s4-burnside-kani-rosen-prym-square}，
\[
\operatorname{Prym}(C_6/H)\sim_{\QQ}J(Z)^2.
\]
故
\[
\dim\operatorname{Prym}(C_6/H)=2g(Z)=8.
\]
若 \(A^2\sim_{\QQ}\operatorname{Prym}(C_6/H)\)，则必有
\[
2\dim A=8,
\qquad
\dim A=4.
\]
另一方面，由结论 \ref{con:xi-time-part9n-s4-all-subgroup-genera} 与结论 \ref{con:xi-time-part9n-s4-hodge-tensor-factorization}，
\[
\dim J(H)=g(H)=5,\qquad
\dim J(Z)=g(Z)=4,
\]
\[
\dim J(C_F)=g(C_F)=3,\qquad
\dim P=9.
\]
在这四个可见因子中，只有 \(J(Z)\) 的维数等于 \(4\)。故 \(A\sim_{\QQ}J(Z)\) 为唯一可能。证毕。
\end{proof}

\endinput

\paragraph{prime-register solenoid 对 Hurwitz 链 Jacobian/Prym 的输入阻断与输出代价}
在局部化数系对偶 \(\Sigma_S=\widehat{G_S}\) 的 Pontryagin 刚性与 \(S_4\)-Hurwitz 商塔的 Jacobian/Prym 同源分解都已确定之后，还可把这两条主线压缩为一组统一的态射结论：正维紧复环面指向有限 prime-register solenoid 直积的输入完全消失，而反向输出则仅按因子个数贡献实相位方向；因此 Hurwitz 链中各 Jacobian/Prym 因子的 prime-register 最小输出代价可被精确求值，并对 Kani--Rosen 同源分解严格可加。

\begin{theorem}[有限 prime-register solenoid 积对正维紧复环面不存在非零输入态射]\label{thm:xi-time-part9n1d-prime-register-input-obstruction}
沿用定理 \ref{thm:killo-localization-circle-dimension-prime-ledger-splitting} 的记号。设
\[
S_1,\dots,S_r\subset \mathbb P
\]
为有限非空素数集，并记
\[
G_{S_i}:=\ZZ[S_i^{-1}],
\qquad
\Sigma_{S_i}:=\widehat{G_{S_i}},
\qquad
K:=\prod_{i=1}^r \Sigma_{S_i}.
\]
则对任意复维 \(g\ge 1\) 的紧复环面 \(A\)，都有
\[
\boxed{\;
\operatorname{Hom}_{\mathrm{cts}}(A,K)=0.\;}
\]
特别地，对任意 Jacobian 或 Prym 都成立。
\end{theorem}
\begin{proof}
由于 \(A\) 是复维 \(g\) 的紧复环面，其 Pontryagin 对偶满足
\[
\widehat A\cong \ZZ^{2g}.
\]
又因为 Pontryagin 对偶把有限直积送为直和，故
\[
\widehat K\cong \bigoplus_{i=1}^r G_{S_i}.
\]
任取连续群同态
\[
f:A\longrightarrow K.
\]
对偶化得到离散群同态
\[
\widehat f:\bigoplus_{i=1}^r G_{S_i}\longrightarrow \ZZ^{2g}.
\]
只需证明任意群同态
\[
\phi:G_{S_i}\longrightarrow \ZZ^{2g}
\]
都恒为零。固定 \(i\) 并取任意 \(p\in S_i\)。对每个 \(n\ge 1\)，由于 \(p^{-n}\in G_{S_i}\)，有
\[
p^n\phi(p^{-n})=\phi(1).
\]
因此
\[
\phi(1)\in \bigcap_{n\ge 1}p^n\ZZ^{2g}=\{0\},
\]
故 \(\phi(1)=0\)。再取任意
\[
x=\frac{a}{b}\in G_{S_i},
\]
其中 \(a\in \ZZ\)，且 \(b\) 的全部素因子都属于 \(S_i\)。则
\[
b\,\phi(x)=\phi(a)=a\,\phi(1)=0.
\]
由于 \(\ZZ^{2g}\) 无挠，只能有 \(\phi(x)=0\)。故 \(\phi=0\)，从而 \(\widehat f=0\)。再对偶化便得 \(f=0\)。证毕。
\end{proof}

\begin{theorem}[有限 prime-register solenoid 积对紧交换 Lie 环面的输出秩与满射阈值]\label{thm:xi-time-part9n1d-prime-register-output-rank-threshold}
\leanverified{paper\_xi\_time\_part9n1d\_prime\_register\_output\_rank\_threshold}
沿用定理 \ref{thm:xi-time-part9n1d-prime-register-input-obstruction} 的记号。设 \(A\) 为实维 \(d\ge 1\) 的紧连通交换 Lie 群。则 \(A\cong \TT^d\)，并且对任意连续群同态
\[
f:K\longrightarrow A
\]
成立：
\begin{enumerate}
  \item \(f(K)\) 是 \(A\) 的一个实维数至多为 \(r\) 的子环面；
  \item 存在连续满射
  \[
  K\twoheadrightarrow A
  \]
  当且仅当 \(r\ge d\)。
\end{enumerate}
等价地，
\[
\boxed{\;
\min\left\{
r\ge 1:\ \exists\ S_1,\dots,S_r\subset \mathbb P\ \text{有限非空},
\prod_{i=1}^r \Sigma_{S_i}\twoheadrightarrow A
\right\}
=d.\;}
\]
对任意复维 \(g\) 的 Jacobian 或 Prym，上式右端即 \(2g\)。
\end{theorem}
\begin{proof}
记
\[
B:=f(K).
\]
每个 \(G_{S_i}\subset \QQ\) 都是无挠群，因此其 Pontryagin 对偶 \(\Sigma_{S_i}\) 连通；故 \(K\) 连通，从而 \(B\) 也是 \(A\) 的连通紧子群。由于 \(A\cong \TT^d\)，可知 \(B\) 必为一个子环面，因而
\[
\widehat B\cong \ZZ^{\dim_{\mathbb R}B}.
\]

由满射
\[
K\twoheadrightarrow B
\]
对偶化得到单射
\[
\widehat B\hookrightarrow \widehat K\cong \bigoplus_{i=1}^r G_{S_i}.
\]
另一方面，定理 \ref{thm:killo-localization-circle-dimension-prime-ledger-splitting} 给出
\[
G_{S_i}\otimes_{\ZZ}\QQ\cong \QQ,
\]
故
\[
\left(\bigoplus_{i=1}^r G_{S_i}\right)\otimes_{\ZZ}\QQ\cong \QQ^r.
\]
因此 \(\widehat B\) 的秩至多为 \(r\)，即
\[
\dim_{\mathbb R}B=\operatorname{rk}_{\ZZ}(\widehat B)\le r.
\]
第 \(1\) 条得证。

若 \(f\) 为满射，则 \(B=A\)，于是
\[
d=\dim_{\mathbb R}A\le r.
\]
反之，若 \(r\ge d\)，则由定理 \ref{thm:killo-localization-circle-dimension-prime-ledger-splitting}，每个 \(\Sigma_{S_i}\) 都带有规范满射
\[
\pi_{S_i}:\Sigma_{S_i}\twoheadrightarrow \TT.
\]
取直积得到
\[
\prod_{i=1}^r\pi_{S_i}:K\twoheadrightarrow \TT^r.
\]
再投影到前 \(d\) 个坐标，即得
\[
K\twoheadrightarrow \TT^d\cong A.
\]
故第 \(2\) 条成立。最后一个等价表述只是把第 \(2\) 条重写为最小因子数公式。证毕。
\end{proof}

\begin{corollary}[Hurwitz 链各 Jacobian/Prym 因子的 prime-register 最小输出代价]\label{cor:xi-time-part9n1d-prime-register-hurwitz-minimal-output-cost}
沿用定理 \ref{thm:killo-s4-genus49-jacobian-computable-isogeny} 与结论 \ref{con:xi-time-part9n-s4-all-subgroup-genera} 的记号。对任意正维紧复环面 \(A\)，定义
\[
r_{\min}(A):=
\min\left\{
r\ge 1:\ \exists\ S_1,\dots,S_r\subset \mathbb P\ \text{有限非空},
\prod_{i=1}^r \Sigma_{S_i}\twoheadrightarrow A
\right\}.
\]
则
\[
\boxed{\;
\begin{aligned}
r_{\min}(J(H))&=10,\qquad r_{\min}(J(C_6))=26,\qquad r_{\min}(J(Z))=8,\\
r_{\min}(J(C_F))&=6,\qquad r_{\min}(P)=18,\qquad r_{\min}(J(X))=98.
\end{aligned}
\;}
\]
\end{corollary}
\begin{proof}
由定理 \ref{thm:xi-time-part9n1d-prime-register-output-rank-threshold}，任意正维紧复环面 \(A\) 的最小 prime-register 输出代价满足
\[
r_{\min}(A)=2\dim_{\mathbb C}A.
\]
而由结论 \ref{con:xi-time-part9n-s4-all-subgroup-genera} 与定理 \ref{thm:killo-s4-genus49-jacobian-computable-isogeny}，
\[
g(H)=5,\qquad g(C_6)=13,\qquad g(Z)=4,\qquad g(C_F)=3,
\]
\[
\dim_{\mathbb C}P=9,\qquad g(X)=49.
\]
对 Jacobian 有 \(\dim_{\mathbb C}J(C)=g(C)\)。将上述维数代入即可得到所示六个精确数值。证毕。
\end{proof}

\begin{theorem}[Kani--Rosen 同源分解上的 prime-register 输出代价严格可加]\label{thm:xi-time-part9n1d-prime-register-output-cost-additivity}
\leanverified{paper\_xi\_time\_part9n1d\_prime\_register\_output\_cost\_additivity}
沿用定理 \ref{thm:killo-s4-genus49-jacobian-computable-isogeny} 的记号。对任意正维紧复环面 \(A\)，定义
\[
\rho(A):=
\min\left\{
r\ge 1:\ \exists\ S_1,\dots,S_r\subset \mathbb P\ \text{有限非空},
\prod_{i=1}^r \Sigma_{S_i}\twoheadrightarrow A
\right\}.
\]
则对同源分解
\[
J(X)\sim_{\QQ}J(H)\times J(Z)^2\times J(C_F)^3\times P^3
\]
有严格可加公式
\[
\boxed{\;
\rho(J(X))
=
\rho(J(H))
+2\,\rho(J(Z))
+3\,\rho(J(C_F))
+3\,\rho(P).\;}
\]
特别地，
\[
\boxed{\;
98=10+2\cdot 8+3\cdot 6+3\cdot 18.\;}
\]
\end{theorem}
\begin{proof}
由定理 \ref{thm:xi-time-part9n1d-prime-register-output-rank-threshold}，任意正维紧复环面 \(A\) 都满足
\[
\rho(A)=2\dim_{\mathbb C}A.
\]
另一方面，由同源分解的维数可加性，
\[
\dim_{\mathbb C}J(X)
=
\dim_{\mathbb C}J(H)
+2\,\dim_{\mathbb C}J(Z)
+3\,\dim_{\mathbb C}J(C_F)
+3\,\dim_{\mathbb C}P.
\]
将上式两边同时乘以 \(2\)，即得 \(\rho\) 的严格可加性。再代入
\[
49=5+2\cdot 4+3\cdot 3+3\cdot 9
\]
便得到数值公式
\[
98=10+2\cdot 8+3\cdot 6+3\cdot 18.
\]
证毕。
\end{proof}

\begin{proposition}[九维 Prym 因子的单寄存器输入消失与单相位输出上界]\label{prop:xi-time-part9n1d-prym-single-register-asymmetry}
沿用定理 \ref{thm:killo-s4-genus49-jacobian-computable-isogeny} 的记号，记
\[
P:=\operatorname{Prym}(Y/Z).
\]
任取有限非空素数集 \(S\subset \mathbb P\)，并记
\[
\Sigma_S:=\widehat{\ZZ[S^{-1}]}.
\]
则有
\[
\boxed{\;
\operatorname{Hom}_{\mathrm{cts}}(P,\Sigma_S)=0,
\qquad
\dim_{\mathbb R}\operatorname{im}(g)\le 1
\quad
(\forall\, g:\Sigma_S\to P).\;}
\]
等价地，单一 prime-register 对 \(P\) 只能输出至多一个实相位方向，而不能从 \(P\) 读取任何非零连续输入。
\end{proposition}
\begin{proof}
由定理 \ref{thm:killo-s4-genus49-jacobian-computable-isogeny}，\(\dim_{\mathbb C}P=9\)。因此把定理 \ref{thm:xi-time-part9n1d-prime-register-input-obstruction} 应用于 \(r=1\) 与 \(A=P\)，立即得到
\[
\operatorname{Hom}_{\mathrm{cts}}(P,\Sigma_S)=0.
\]
再把定理 \ref{thm:xi-time-part9n1d-prime-register-output-rank-threshold} 应用于 \(r=1\) 与 \(A=P\)，便得任意连续群同态
\[
g:\Sigma_S\to P
\]
的像都是 \(P\) 中实维数至多为 \(1\) 的子环面。证毕。
\end{proof}

\noindent
上述五条结果表明，局部化数系 prime-register solenoid 与 \(S_4\)-Hurwitz 链 Jacobian/Prym 的接口在输入与输出两端呈现严格非对称性：输入方向不存在任何非零连续态射，而输出方向中每个 prime-register 恰只贡献一个实相位。因而 prime-register 最小输出代价不仅对各个 Jacobian/Prym 因子可精确求值，而且与 Kani--Rosen 同源分解逐块严格兼容。

\endinput


\noindent
结合结论 \ref{con:xi-time-part9n-s4-all-subgroup-genera} 与以上诸条结论，\(S_4\)-Hurwitz 商塔中这条以 Chevalley--Weil 重数为起点的链条已在几何、表示、算术与 prime-register 态射预算四侧同时闭合：一方面，\(C_6\) 作为 \(H\) 与 \(Z\) 的规范化纤维积被唯一识别，全纯微分空间按四个非平凡不可约表示张量分层，八个商曲线的 genus 以及 \(X/C_2\)、\(X/C_3\) 等此前缺失商的 Jacobian 同源型都由 Burnside--Kani--Rosen 关系与同一组表示论数据唯一补全，从而四个非平凡不可约表示的几何承载因子不再遗留任何隐藏等型块；另一方面，\(J(X)\) 的有理投影代数下界、\(\operatorname{Prym}(C_6/H)\) 的平方根刚性、良素处局部 Euler 因子的平方/立方块、两类 Prym 的最小同源代价及其 \(3\)-初等与 \(2\)-初等极化缺陷，以及 \(V_3'\)-等型三重块 \(P^3\) 的主极化化次数、\(\QQ\)-子等型分类与 Rosati 对称端代数，继续由群论分块与极化缺陷严格锁定；此外，有限 prime-register solenoid 对正维 Jacobian/Prym 的输入态射完全消失，而其输出代价恰等于目标实维数，并对 Kani--Rosen 同源分解逐块保持严格可加。

\endinput


\begin{conclusion}[\texorpdfstring{$\operatorname{Prym}(C_6/H)$}{Prym(C6/H)} 的 \texorpdfstring{$\QQ(\zeta_3)$}{Q(zeta3)}-Weil 型结构与签名 \texorpdfstring{$(4,4)$}{(4,4)}]\label{con:xi-time-part9n-prym-c6h-weil-type-signature}
阿贝尔簇 \(A\) 是以
\[
K:=\QQ(\zeta_3)
\]
为虚二次域的 Weil 型阿贝尔八维簇，其 \(K\)-signature 为
\[
(4,4).
\]
并且 \(A\) 的诱导极化型为
\[
(1,1,1,1,3,3,3,3),
\]
且有
\[
\QQ(\zeta_3)\subset \operatorname{End}^0(A),\qquad
A\sim_{\QQ}J(Z)^2.
\]
特别地，存在自然二维 \(\QQ\)-向量空间
\[
W_K:=\bigwedge\nolimits_{K}^{8}H^1(A,\QQ)\subset H^8(A,\QQ)\cap H^{4,4}(A),
\]
其由 Weil classes 给出；经同源 \(A\sim_{\QQ}J(Z)^2\)，该二维 Hodge 类空间亦传递到
\[
H^8\!\bigl(J(Z)^2,\QQ\bigr)\cap H^{4,4}\!\bigl(J(Z)^2\bigr).
\]
\end{conclusion}
\begin{proof}
由定理 \ref{thm:killo-s4-even-subgroup-free-action-unramified-a4-lift}，覆盖
\[
\pi:C_6\to H
\]
为无分歧循环三重覆盖，且
\[
g(H)=5,\qquad g(C_6)=13.
\]
故
\[
\dim A=g(C_6)-g(H)=8.
\]
再由定理 \ref{thm:killo-s4-burnside-kani-rosen-prym-square} 与命题 \ref{prop:killo-two-prym-polarization-types}，
\[
A\sim_{\QQ}J(Z)^2,
\]
并且 \(A\) 的诱导极化型为 \((1^4,3^4)\)，同时
\[
\QQ(\zeta_3)\subset \operatorname{End}^0(A).
\]

为确定 signature，记 \(\sigma\) 为 \(\pi\) 的甲板变换生成元。由于 \(\pi\) 为循环三重无分歧覆盖，存在非常值 \(3\)-扭线丛
\[
L\in \operatorname{Pic}^0(H),\qquad L^{\otimes 3}\simeq \mathcal O_H,
\]
使得
\[
\pi_*\mathcal O_{C_6}\simeq \mathcal O_H\oplus L^{-1}\oplus L^{-2}.
\]
由投影公式，
\[
\pi_*\omega_{C_6}
\simeq
\omega_H\oplus (\omega_H\otimes L)\oplus (\omega_H\otimes L^2).
\]
于是 \(H^0(C_6,\Omega_{C_6}^1)\) 按 \(C_3=\langle \sigma\rangle\) 的特征分解为
\[
H^0(C_6,\Omega_{C_6}^1)=U_0\oplus U_{\chi}\oplus U_{\bar\chi},
\qquad
\chi(\sigma)=\zeta_3,
\]
其中
\[
U_0\simeq H^0(H,\omega_H),\qquad
U_{\chi}\simeq H^0(H,\omega_H\otimes L),\qquad
U_{\bar\chi}\simeq H^0(H,\omega_H\otimes L^2).
\]
由于 \(L\) 与 \(L^2\) 均为非平凡零次数线丛，Riemann--Roch 与 Serre 对偶给出
\[
h^0(H,\omega_H\otimes L)=h^0(H,\omega_H\otimes L^2)=g(H)-1=4,
\]
而
\[
h^0(H,\omega_H)=g(H)=5.
\]
故
\[
\dim U_0=5,\qquad \dim U_{\chi}=\dim U_{\bar\chi}=4.
\]
Prym 余切空间正对应于非平凡特征部分
\[
H^0(A,\Omega_A^1)\simeq U_{\chi}\oplus U_{\bar\chi},
\]
因此 \(K\)-作用在 \(H^{1,0}(A)\) 上的两条复嵌入重数均为 \(4\)，即 signature 为 \((4,4)\)。

另一方面，甲板变换群 \(C_3\) 通过保持 Prym 诱导极化的自同构作用于 \(A\)，故 Rosati 对合把 \(\sigma\) 送到 \(\sigma^{-1}\)，也即在 \(K=\QQ(\zeta_3)\) 上对应复共轭。于是 \(A\) 在该极化下确为 \(K\)-Weil 型阿贝尔八维簇。对签名 \((4,4)\) 的 Weil 型阿贝尔簇，一般理论表明
\[
W_K=\bigwedge\nolimits_K^{8}H^1(A,\QQ)
\]
是二维 \(\QQ\)-向量空间，并落在 \(H^8(A,\QQ)\cap H^{4,4}(A)\) 中。最后，同源 \(A\sim_{\QQ}J(Z)^2\) 在有理上同调上诱导 Hodge 结构同构，故 \(W_K\) 传递到 \(H^8(J(Z)^2,\QQ)\) 中相应的 \((4,4)\)-部分。证毕。
\end{proof}

\begin{conclusion}[两类 Prym 缺陷核在素数支撑上的完全分离]\label{con:xi-time-part9n-prym-polarization-primary-splitting}
记 \(\lambda_A\) 与 \(\lambda_P\) 分别为命题 \ref{prop:killo-two-prym-polarization-types} 在
\[
A=\operatorname{Prym}(C_6/H),\qquad P=\operatorname{Prym}(Y/Z)
\]
上给出的诱导极化。则有
\[
\ker(\lambda_A)\cong (\ZZ/3\ZZ)^8,
\qquad
\ker(\lambda_P)\cong (\ZZ/2\ZZ)^{10}.
\]
因此，在可计算同源分解
\[
J(X)\sim_{\QQ}J(H)\times J(Z)^2\times J(C_F)^3\times P^3
\]
中，显式出现的非主极化缺陷按素数严格解耦：\(3\)-初缺陷全部锚定在
\[
A\sim_{\QQ}J(Z)^2
\]
这一 Prym 平方方向，而 \(2\)-初缺陷全部锚定在 \(P\) 方向；相对地，\(J(H)\) 与 \(J(C_F)\) 作为 Jacobian 因子各自携带其自然的主极化。
\end{conclusion}
\begin{proof}
对任意极化型
\[
(d_1,\dots,d_g),
\]
其极化同态的核满足标准结构公式
\[
\ker(\lambda)\cong \bigoplus_{i=1}^{g}(\ZZ/d_i\ZZ)^2.
\]
由命题 \ref{prop:killo-two-prym-polarization-types}，
\[
\mathrm{type}(\lambda_A)=(1,1,1,1,3,3,3,3),
\]
故
\[
\ker(\lambda_A)\cong (\ZZ/3\ZZ)^8.
\]
同理，
\[
\mathrm{type}(\lambda_P)=(1,1,1,1,2,2,2,2,2),
\]
从而
\[
\ker(\lambda_P)\cong (\ZZ/2\ZZ)^{10}.
\]
再由定理 \ref{thm:killo-s4-burnside-kani-rosen-prym-square} 与
\ref{thm:killo-s4-genus49-jacobian-computable-isogeny}，
\[
A\sim_{\QQ}J(Z)^2,\qquad
J(X)\sim_{\QQ}J(H)\times J(Z)^2\times J(C_F)^3\times P^3.
\]
故与 \(A\) 对应的缺陷群是纯 \(3\)-群，而与 \(P\) 对应的缺陷群是纯 \(2\)-群，二者在素数层面彼此互素。最后，\(J(H)\) 与 \(J(C_F)\) 作为 Jacobian，本身都带有其规范主极化，故不引入新的极化缺陷素因子。证毕。
\end{proof}

\begin{conclusion}[Burnside--Kani--Rosen 同源关系诱导的局部 \texorpdfstring{$\zeta$}{zeta} 因子分解与点数恒等式]\label{con:xi-time-part9n-burnside-kani-rosen-local-zeta-pointcount}
设 \(v\) 为一处使
\[
X,\ H,\ C_6,\ Z,\ C_F,\ Y,\ X/C_2,\ X/C_3
\]
及阿贝尔簇 \(P\) 同时良约化的有限素位，记残域大小为 \(q_v\)。对任意光滑射影曲线 \(C/\FF_{q_v}\)，定义
\[
Z(C/\FF_{q_v},T)=\frac{P_C(v,T)}{(1-T)(1-q_vT)},
\]
其中 \(P_C(v,T)\) 为 Frobenius 在 \(H^1_{\text{ét}}(C_{\overline{\FF}_{q_v}},\QQ_\ell)\) 上的特征多项式；并记
\[
P_P(v,T):=
\det\!\Bigl(1-T\,\mathrm{Frob}_v\mid
H^1_{\text{ét}}(P_{\overline{\FF}_{q_v}},\QQ_\ell)\Bigr).
\]
则有
\[
P_{C_6}(v,T)=P_H(v,T)\,P_Z(v,T)^2,
\]
\[
P_{X/C_2}(v,T)=P_{C_F}(v,T)^2\,P_Y(v,T),
\]
\[
P_{C_6}(v,T)\,P_{X/C_3}(v,T)^2
=
P_H(v,T)^3\,P_{C_F}(v,T)^2\,P_Y(v,T)^2,
\]
以及
\[
P_X(v,T)=P_H(v,T)\,P_Z(v,T)^2\,P_{C_F}(v,T)^3\,P_P(v,T)^3.
\]

等价地，对任意 \(n\ge 1\)，若记
\[
N_C(v,n):=\#C(\FF_{q_v^n}),
\qquad
a_C(v,n):=q_v^n+1-N_C(v,n),
\]
则有精确恒等式
\[
a_{C_6}(v,n)=a_H(v,n)+2a_Z(v,n),
\]
\[
a_{X/C_2}(v,n)=2a_{C_F}(v,n)+a_Y(v,n),
\]
\[
a_{C_6}(v,n)+2a_{X/C_3}(v,n)
=
3a_H(v,n)+2a_{C_F}(v,n)+2a_Y(v,n).
\]
\end{conclusion}
\begin{proof}
由定理 \ref{thm:killo-s4-burnside-kani-rosen-prym-square}，
\[
J(C_6)\sim_{\QQ}J(H)\times J(Z)^2,
\]
\[
J(X/C_2)\sim_{\QQ}J(C_F)^2\times J(Y),
\]
\[
J(C_6)\times J(X/C_3)^2
\sim_{\QQ}
J(H)^3\times J(C_F)^2\times J(Y)^2.
\]
再由定理 \ref{thm:killo-s4-genus49-jacobian-computable-isogeny}，
\[
J(X)\sim_{\QQ}J(H)\times J(Z)^2\times J(C_F)^3\times P^3.
\]
在共同良约化处，同源阿贝尔簇的 \(\ell\)-进 \(H^1\) 具有相同的 Frobenius 特征多项式，而直积对应特征多项式相乘，故立即得到四条 \(P\)-因子恒等式。

对任意曲线 \(C/\FF_{q_v}\)，标准展开给出
\[
\log P_C(v,T)=-\sum_{n\ge 1}\frac{a_C(v,n)}{n}T^n.
\]
将前三条 \(P\)-因子恒等式取对数并比较 \(T^n\) 系数，即得三条点数恒等式。证毕。
\end{proof}

\begin{conclusion}[九维 Prym 因子的 Frobenius 迹整除律]\label{con:xi-time-part9n-prym9-frobenius-trace-divisibility}
对任意共同良约化素位 \(v\) 及任意 \(n\ge 1\)，有严格公式
\[
\operatorname{Tr}\!\Bigl(
\mathrm{Frob}_v^n\mid
H^1_{\text{ét}}(P_{\overline{\FF}_{q_v}},\QQ_\ell)
\Bigr)
=
\frac{
a_X(v,n)-a_H(v,n)-2a_Z(v,n)-3a_{C_F}(v,n)
}{3}.
\]
特别地，
\[
a_X(v,n)-a_H(v,n)-2a_Z(v,n)-3a_{C_F}(v,n)\equiv 0\pmod 3.
\]
\end{conclusion}
\begin{proof}
由定理 \ref{thm:killo-s4-genus49-jacobian-computable-isogeny}，
\[
J(X)\sim_{\QQ}J(H)\times J(Z)^2\times J(C_F)^3\times P^3.
\]
因此在共同良约化处，
\[
\operatorname{Tr}\!\Bigl(
\mathrm{Frob}_v^n\mid H^1_{\text{ét}}(J(X)_{\overline{\FF}_{q_v}},\QQ_\ell)
\Bigr)
=
a_H(v,n)+2a_Z(v,n)+3a_{C_F}(v,n)
\]
\[
\qquad\qquad
+3\,\operatorname{Tr}\!\Bigl(
\mathrm{Frob}_v^n\mid H^1_{\text{ét}}(P_{\overline{\FF}_{q_v}},\QQ_\ell)
\Bigr).
\]
又对曲线 Jacobian 有
\[
H^1_{\text{ét}}(J(X)_{\overline{\FF}_{q_v}},\QQ_\ell)
\cong
H^1_{\text{ét}}(X_{\overline{\FF}_{q_v}},\QQ_\ell),
\]
故左端正是 \(a_X(v,n)\)。移项并除以 \(3\) 即得所示公式。右端既等于实际的 Frobenius 迹，因而必为整数，从而得到模 \(3\) 整除律。证毕。
\end{proof}

\begin{conclusion}[属 \(49\) Jacobian 上由 \texorpdfstring{$\QQ(\zeta_3)$}{Q(zeta3)}-Weil 型强制的余维 \texorpdfstring{$4$}{4} Hodge 类]\label{con:xi-time-part9n-jx-codim4-weil-hodge-classes}
在 \(H^8(J(X),\QQ)\) 中存在自然二维 \(\QQ\)-子空间
\[
W_{K,X}\subset H^8(J(X),\QQ)\cap H^{4,4}(J(X)),
\qquad K=\QQ(\zeta_3),
\]
它由结论 \ref{con:xi-time-part9n-prym-c6h-weil-type-signature} 中 \(A=\operatorname{Prym}(C_6/H)\) 的 Weil classes 经同源分解
\[
J(X)\sim_{\QQ}J(H)\times J(Z)^2\times J(C_F)^3\times P^3
\]
传递而来。换言之，属 \(49\) 的 Jacobian \(J(X)\) 含有一组由无分歧循环三重覆盖 \(C_6\to H\) 强制产生的余维 \(4\) Hodge 类。
\end{conclusion}
\begin{proof}
由结论 \ref{con:xi-time-part9n-prym-c6h-weil-type-signature}，存在二维 \(\QQ\)-向量空间
\[
W_K\subset H^8(A,\QQ)\cap H^{4,4}(A),
\qquad A=\operatorname{Prym}(C_6/H).
\]
同一结论与定理 \ref{thm:killo-s4-burnside-kani-rosen-prym-square} 给出
\[
A\sim_{\QQ}J(Z)^2.
\]
故 \(W_K\) 可视为
\[
H^8\!\bigl(J(Z)^2,\QQ\bigr)\cap H^{4,4}\!\bigl(J(Z)^2\bigr)
\]
中的二维 Hodge 子结构。

另一方面，定理 \ref{thm:killo-s4-genus49-jacobian-computable-isogeny} 给出
\[
J(X)\sim_{\QQ}J(H)\times J(Z)^2\times J(C_F)^3\times P^3.
\]
该同源分解在有理上同调上诱导 Hodge 结构同构，因此
\[
H^8\!\bigl(J(Z)^2,\QQ\bigr)
\]
自然识别为 \(H^8(J(X),\QQ)\) 的一个有理 Hodge 直和因子。将上述 \(W_K\) 沿此同构拉回，即得
\[
W_{K,X}\subset H^8(J(X),\QQ)\cap H^{4,4}(J(X)),
\]
且 \(\dim_{\QQ}W_{K,X}=2\)。证毕。
\end{proof}

同一组 \(S_4\)-Hurwitz 商塔数据还进一步固定了权一表示的等型重数、若干未显式写出的商曲线亏格、八维 Prym 的端同态约束以及 \(J(C_6)\) 在 Siegel 模空间中的位置；相应地可压缩为下列四条结论。

\begin{conclusion}[\texorpdfstring{$S_4$}{S4} 权一表示的完全分解与全部商曲线亏格的统一恢复]\label{con:xi-time-part9n-s4-hodge-quotient-genera-recovery}
令
\[
V:=H^0(X,\Omega_X^1).
\]
则有有理 \(S_4\)-表示的完全分解
\[
\boxed{\;
V\cong
\mathrm{sgn}^{\oplus 5}\oplus
V_2^{\oplus 4}\oplus
V_3^{\oplus 3}\oplus
\bigl(V_3'\bigr)^{\oplus 9}.\;}
\]
并且对任意子群 \(K\le S_4\)，若记
\[
\rho_K:=\Ind_K^{S_4}\mathbf 1,
\]
则对应商曲线亏格满足
\[
g(X/K)=\dim V^K=\langle V,\rho_K\rangle_{S_4}.
\]
特别地，
\[
\boxed{\;
g(X/C_2)=19,\qquad g(X/C_3)=17.\;}
\]
与此同时，文中已出现的四个商曲线亏格
\[
g(H)=5,\qquad g(Z)=4,\qquad g(C_F)=3,\qquad g(Y)=13
\]
亦由同一角色公式统一恢复。
\end{conclusion}
\begin{proof}
把定理 \ref{thm:killo-s4-genus49-jacobian-computable-isogeny} 所用的等型分块记为
\[
J(X)\sim A_{\mathrm{sgn}}\times A_2\times A_3\times A_{3'}.
\]
由同一定理及其证明中的商曲线识别，
\[
A_{\mathrm{sgn}}\sim J(H),\qquad
A_2\sim J(Z)^2,\qquad
A_3\sim J(C_F)^3,\qquad
A_{3'}\sim P^3.
\]
因此
\[
\dim A_{\mathrm{sgn}}=g(H)=5,\qquad
\dim A_2=2g(Z)=8,
\]
\[
\dim A_3=3g(C_F)=9,\qquad
\dim A_{3'}=3\dim P=27.
\]
若 \(m_\rho\) 表示不可约有理表示 \(\rho\) 在 \(V=H^0(X,\Omega_X^1)\) 中的重数，则对应等型阿贝尔因子的维数满足
\[
\dim A_\rho=(\dim \rho)\,m_\rho.
\]
于是
\[
m_{\mathrm{sgn}}=5,\qquad
m_{V_2}=4,\qquad
m_{V_3}=3,\qquad
m_{V_3'}=9.
\]
又因 \(X/S_4\simeq \PP^1\)，故平凡表示重数为 \(0\)。遂得所示完全分解。

对任意 \(K\le S_4\)，Frobenius 互反给出
\[
\dim V^K=\langle V,\rho_K\rangle_{S_4},
\qquad
\rho_K=\Ind_K^{S_4}\mathbf 1.
\]
将
\[
\rho_{A_4}=\mathbf 1\oplus \mathrm{sgn},\qquad
\rho_{D_8}=\mathbf 1\oplus V_2,\qquad
\rho_{S_3}=\mathbf 1\oplus V_3,
\]
\[
\rho_{C_4}=\mathbf 1\oplus V_2\oplus V_3',\qquad
\rho_{C_2}=\mathbf 1\oplus V_2\oplus 2V_3\oplus V_3',
\]
\[
\rho_{C_3}=\mathbf 1\oplus \mathrm{sgn}\oplus V_3\oplus V_3'
\]
代入即得
\[
g(H)=5,\qquad g(Z)=4,\qquad g(C_F)=3,\qquad g(Y)=4+9=13,
\]
并且
\[
g(X/C_2)=4+2\cdot 3+9=19,
\qquad
g(X/C_3)=5+3+9=17.
\]
证毕。
\end{proof}

\begin{conclusion}[\texorpdfstring{$\operatorname{Prym}(C_6/H)$}{Prym(C6/H)} 的最小端同态表述与平方同源不相容]\label{con:xi-time-part9n-prym-c6h-endomorphism-obstruction}
若将猜想 \ref{conj:killo-minimal-extra-endomorphism-large-prym} 的第二条按字面解释为
\[
\operatorname{End}^0_{\overline{\QQ}}\!\bigl(\operatorname{Prym}(C_6/H)\bigr)=\QQ(\zeta_3),
\]
则该表述与定理 \ref{thm:killo-s4-burnside-kani-rosen-prym-square} 不相容。更精确地，
\[
\boxed{\;
\dim_{\QQ}
\operatorname{End}^0_{\overline{\QQ}}\!\bigl(\operatorname{Prym}(C_6/H)\bigr)
\ge 4.\;}
\]
因此，八维 Prym 的“最小额外自同态”不可能直接在整个端同态代数层面收缩为二次数域 \(\QQ(\zeta_3)\)。
\end{conclusion}
\begin{proof}
由定理 \ref{thm:killo-s4-burnside-kani-rosen-prym-square}，
\[
\operatorname{Prym}(C_6/H)\sim_{\QQ} J(Z)^2.
\]
同源阿贝尔簇的有理端同态代数经由准同源传递同构，因此
\[
\operatorname{End}^0_{\overline{\QQ}}\!\bigl(\operatorname{Prym}(C_6/H)\bigr)
\cong
\operatorname{End}^0_{\overline{\QQ}}\!\bigl(J(Z)^2\bigr).
\]
而对任意正维阿贝尔簇 \(B\)，都有标准恒等式
\[
\operatorname{End}^0_{\overline{\QQ}}(B^2)\cong
M_2\!\bigl(\operatorname{End}^0_{\overline{\QQ}}(B)\bigr).
\]
取 \(B=J(Z)\)，即得
\[
\operatorname{End}^0_{\overline{\QQ}}\!\bigl(\operatorname{Prym}(C_6/H)\bigr)
\cong
M_2\!\bigl(\operatorname{End}^0_{\overline{\QQ}}(J(Z))\bigr).
\]
由于恒等元给出
\[
\dim_{\QQ}\operatorname{End}^0_{\overline{\QQ}}(J(Z))\ge 1,
\]
故
\[
\dim_{\QQ}
\operatorname{End}^0_{\overline{\QQ}}\!\bigl(\operatorname{Prym}(C_6/H)\bigr)
=
4\,
\dim_{\QQ}\operatorname{End}^0_{\overline{\QQ}}(J(Z))
\ge 4.
\]
另一方面，
\[
\dim_{\QQ}\QQ(\zeta_3)=2.
\]
因此不可能有
\[
\operatorname{End}^0_{\overline{\QQ}}\!\bigl(\operatorname{Prym}(C_6/H)\bigr)=\QQ(\zeta_3).
\]
证毕。
\end{proof}

\begin{conclusion}[\texorpdfstring{$V_2$}{V2} 等型块的平衡 \texorpdfstring{$\QQ(\zeta_3)$}{Q(zeta3)}-Weil 结构]\label{con:xi-time-part9n-v2-block-balanced-weil}
令 \(A_2\) 表示 \(J(X)\) 的 \(V_2\)-等型块。则
\[
A_2\sim_{\QQ}J(Z)^2\sim_{\QQ}\operatorname{Prym}(C_6/H).
\]
并且经由同源识别
\[
A_2\sim_{\QQ}\operatorname{Prym}(C_6/H)
\]
所传递的 \(C_3\)-作用满足
\[
H^0(A_2,\Omega_{A_2}^1)\cong \chi^{\oplus 4}\oplus \bar\chi^{\oplus 4},
\]
其中 \(\chi,\bar\chi\) 为 \(C_3\) 的两个非平凡复特征。因而 \(A_2\) 是一个以
\[
K=\QQ(\zeta_3)
\]
为乘法域的八维 Weil 型阿贝尔簇，其 \(K\)-signature 恰为
\[
\boxed{\;(4,4).\;}
\]
此外，\(A_2\) 的诱导极化型为
\[
(1,1,1,1,3,3,3,3).
\]
\end{conclusion}
\begin{proof}
结论 \ref{con:xi-time-part9n-prym-c6h-weil-type-signature} 已证明
\[
\operatorname{Prym}(C_6/H)
\]
是 \(K=\QQ(\zeta_3)\)-Weil 型阿贝尔八维簇，且其签名为 \((4,4)\)，并具有极化型 \((1^4,3^4)\)。另一方面，定理 \ref{thm:killo-s4-burnside-kani-rosen-prym-square} 给出
\[
\operatorname{Prym}(C_6/H)\sim_{\QQ}J(Z)^2,
\]
而定理 \ref{thm:killo-s4-genus49-jacobian-computable-isogeny} 正是将 \(J(Z)^2\) 识别为 \(J(X)\) 的 \(V_2\)-等型块 \(A_2\)。故所述 Weil 结构、特征分解与极化型全部传递到 \(A_2\)。证毕。
\end{proof}

\begin{conclusion}[\texorpdfstring{$J(C_6)$}{J(C6)} 落在余维至少 \texorpdfstring{$11$}{11} 的特异分解模子簇中]\label{con:xi-time-part9n-c6-special-jacobian-locus}
设
\[
\mathcal D_{5,(4)^2}
:=
\left\{
A\in \mathcal A_{13}:\ 
A\sim B_5\times B_4^2
\text{ for some }B_5\in\mathcal A_5,\ B_4\in\mathcal A_4
\right\}.
\]
则
\[
J(C_6)\in \mathcal D_{5,(4)^2},
\]
并且
\[
\dim \mathcal D_{5,(4)^2}\le 25.
\]
因此，\(J(C_6)\) 在属 \(13\) 的 Jacobian 模子簇
\[
\mathcal J_{13}\subset \mathcal A_{13}
\]
中落在余维至少
\[
\boxed{\;36-25=11\;}
\]
的特异子簇内。换言之，\(C_6\) 并非属 \(13\) 的泛性曲线，而被强制压入一个高度特异的同源分解子簇。
\end{conclusion}
\begin{proof}
由定理 \ref{thm:killo-s4-burnside-kani-rosen-prym-square}，
\[
J(C_6)\sim J(H)\times J(Z)^2,
\]
其中 \(g(H)=5\)、\(g(Z)=4\)。故 \(J(C_6)\) 的确属于 \(\mathcal D_{5,(4)^2}\)。

再看维数。Siegel 模空间满足
\[
\dim \mathcal A_g=\frac{g(g+1)}{2},
\]
故
\[
\dim \mathcal A_5=15,\qquad \dim \mathcal A_4=10.
\]
映射
\[
\mathcal A_5\times \mathcal A_4\longrightarrow \mathcal A_{13},
\qquad
(B_5,B_4)\longmapsto B_5\times B_4^2
\]
的像维数至多为 \(15+10=25\)。允许“\(\sim\)”意义下的同源，只会产生该乘积像的可数个 Hecke 对应平移，因此相应同源子簇的维数仍至多为 \(25\)。于是
\[
\dim \mathcal D_{5,(4)^2}\le 25.
\]

另一方面，属 \(13\) 的 Jacobian 子簇 \(\mathcal J_{13}\) 与曲线模空间 \(\mathcal M_{13}\) 具有相同维数
\[
\dim \mathcal J_{13}=\dim \mathcal M_{13}=3\cdot 13-3=36.
\]
故 \(J(C_6)\) 所在的分解子簇在 \(\mathcal J_{13}\) 中的余维至少为
\[
36-25=11.
\]
证毕。
\end{proof}
